% Mechanically extracted after \mainmatter from the frozen standalone master.

\markboth{Exposé XXIII. Uniqueness of Pinned Groups}
         {Exposé XXIII. Uniqueness of Pinned Groups}
% BEGIN inputs/XXIII/tex/components/00_expose_XXIII_title_and_definition11.tex
\chapter*{Exposé XXIII\\Reductive Groups: Uniqueness of Pinned Groups}
\addcontentsline{toc}{chapter}{Exposé XXIII. Uniqueness of Pinned Groups}
\begin{center}
\textit{by M. Demazure}
\end{center}

\NDE{Version of 13 October 2024.}
The aim of this Exposé is to prove the uniqueness theorem
(Theorem~4.1). It was proved by Chevalley in the case of an
algebraically closed field; the reduction-to-rank-two method used here
is likewise due to Chevalley (see \textit{Bible}, Exposés~23 and~24).
Along the way, we obtain an explicit description of reductive groups
by generators and relations (3.5).

\SGANextSectionTarget{sga3:xxiii:section:1}\SGASection{1}{Pinnings}

\SGARefTarget{sga3:XXIII:1.1}
\SGAStatement{Definition 1.1}{Let \(S\) be a scheme and let
\((G,T,M,R)\) be an \(S\)-split group (XXII, \SGARefLink{sga3:xxii:def:1.13}{1.13}). A \emph{pinning}
\NDE{Demazure tells us that behind this terminology lies the image of
the butterfly (supplied to him by Grothendieck): the body is a maximal
torus \(T\), the wings are two Borel subgroups opposite with respect to
\(T\); one spreads the butterfly by laying out the wings and then fixes
elements in the additive groups (pins) to rigidify the situation,
that is, to eliminate the automorphisms.}
of this split group is the datum of a system \(\Delta\) of simple
roots of \(R\), together with, for each \(\alpha\in\Delta\), a section
\[
  X_\alpha\in\Gamma(S,\mathfrak g^\alpha)^\times .
\]}

In other words, a pinning of the reductive group \(G\) over the
nonempty scheme \(S\) is the datum of:
\begin{enumerate}
  \item[(i)] a maximal torus \(T\);
  \item[(ii)] an abelian group \(M\) and an isomorphism
    \(T\simeq D_S(M)\);
  \item[(iii)] a system of roots \(R\) of \(G\) with respect to \(T\);
  \item[(iv)] a system \(\Delta\) of simple roots of \(R\);
  \item[\SGAInlineRefTarget{sga3:XXIII:1.1:item:v}(v)] an element
    \(X_\alpha\in\Gamma(S,\mathfrak g^\alpha)^\times\) for every
    \(\alpha\in\Delta\), that is, an element
    \[
      u_\alpha=\exp_\alpha(X_\alpha)\in U_\alpha^\times(S)
      \qquad\text{for every }\alpha\in\Delta,
    \]
\end{enumerate}
satisfying condition (D1) of \SGARefLink{sga3:xxii:def:1.13}{Exposé~XXII, 1.13} (indeed, condition
(D2) of \textit{loc.\ cit.} is automatic\NDE{It is implied by
condition \SGARefLink{sga3:XXIII:1.1:item:v}{(v)}, namely the existence of a section
\(X_\alpha\in\Gamma(S,\mathfrak g^\alpha)^\times\).}).

Every split group admits a pinning; in particular, every reductive
group is locally pinnable for the étale topology.

% END inputs/XXIII/tex/components/00_expose_XXIII_title_and_definition11.tex

% BEGIN inputs/XXIII/tex/components/01_expose_XXIII_par12_and_definition13.tex
\SGARefTarget{sga3:XXIII:1.2}
\SGAParagraph{1.2}
If \(G\) is a pinned \(S\)-group, that is, an \(S\)-split group
equipped with a pinning, it is canonically endowed with the system of
positive roots \(R^+\) defined by \(\Delta\), the corresponding Borel
subgroup \(B=B_{R^+}\), the opposite Borel subgroup \(B^-=B_{R^-}\),
the unipotent groups \(U=B_u\) and \(U^-=(B^-)_{u}\), the open
subscheme \(U^-\cdot T\cdot U\), and so forth. Likewise, for each
\(\alpha\in\Delta\), there is a canonical isomorphism of vector groups
\[
  p_\alpha:\mathbb G_{a,S}\xrightarrow{\ \sim\ }U_\alpha,
  \qquad x\longmapsto \exp_\alpha(xX_\alpha)=u_\alpha^x,
\]
normalized by \(T\) with multiplier \(\alpha\), and whose datum is
equivalent to that of \(X_\alpha\) (\SGARefLink{sga3:xxii:thm:1.1}{Exposé~XXII, 1.1}).

By duality one obtains an
\(X_{-\alpha}\in\Gamma(S,\mathfrak g^{-\alpha})^\times\) and an
isomorphism
\[
  p_{-\alpha}:\mathbb G_{a,S}\xrightarrow{\ \sim\ }U_{-\alpha}
\]
which is contragredient to the preceding one (\SGARefLink{sga3:xxii:not:1.3}{Exposé~XXII, 1.3}). We
set (\SGARefLink{sga3:xx:3.1}{Exposé~XX, 3.1})
\[
  w_\alpha=w_\alpha(X_\alpha)
  =p_\alpha(1)p_{-\alpha}(-1)p_\alpha(1)
  =p_{-\alpha}(-1)p_\alpha(1)p_{-\alpha}(-1).
\]
Then (\textit{loc.\ cit.}, \SGARefLink{sga3:xxiii:subsection:3-1}{3.1} and \SGARefLink{sga3:xx:3.7}{3.7})
\[
  w_\alpha^2=\alpha^\ast(-1),\qquad
  \operatorname{int}(w_\alpha)t
    =s_\alpha(t)=t\cdot\alpha^\ast\!\bigl(\alpha(t)^{-1}\bigr),
\]
\[
\begin{aligned}
  \operatorname{int}(w_\alpha)p_\alpha(x)
    &=p_{-\alpha}(-x)=p_{-\alpha}(x)^{-1},&
  \operatorname{Ad}(w_\alpha)X_\alpha&=-X_{-\alpha},\\
  \operatorname{int}(w_\alpha)p_{-\alpha}(x)
    &=p_\alpha(-x)=p_\alpha(x)^{-1},&
  \operatorname{Ad}(w_\alpha)X_{-\alpha}&=-X_\alpha .
\end{aligned}
\]
We shall systematically use the preceding notation in what follows.

\SGARefTarget{sga3:XXIII:1.3}
\SGAStatement{Definition 1.3}{Let \(S\) be a scheme, and suppose that
\[
  (G,T,M,R,\Delta,(X_\alpha))
  \quad\text{and}\quad
  (G',T',M',R',\Delta',(X'_{\alpha'}))
\]
are two pinned \(S\)-groups.
The morphism of split groups (\SGARefLink{sga3:xxii:def:4.2.1}{Exposé~XXII, 4.2.1})
\[
  f:G\longrightarrow G'
\]
is said to be \emph{compatible with the pinnings}, or to define a
\emph{morphism of pinned groups}, if the associated bijection
\(d:R\to R'\) (cf.\ \textit{loc.\ cit.}) satisfies
\(d(\Delta)=\Delta'\) and if, for every \(\alpha\in\Delta\),
\[
  f\bigl(\exp_\alpha(X_\alpha)\bigr)
    =\exp_{d(\alpha)}\bigl(X'_{d(\alpha)}\bigr),
  \qquad\text{i.e.}\qquad
  f(u_\alpha)=u'_{d(\alpha)} .
\]}
\NDE{The bijection \(R\xrightarrow{\sim}R'\) is denoted by \(d\)
(rather than by \(u\)) in order to avoid the notation
\(u'_{u(\alpha)}\).}

% END inputs/XXIII/tex/components/01_expose_XXIII_par12_and_definition13.tex

% BEGIN inputs/XXIII/tex/components/02_expose_XXIII_par14_through_scholium16.tex
\SGARefTarget{sga3:XXIII:1.4}
\SGAParagraph{1.4}
If \(q(\alpha)\) denotes the integer of \textit{loc.\ cit.}, then
\[
  f\bigl(p_\alpha(x)\bigr)
    =p'_{d(\alpha)}\bigl(x^{q(\alpha)}\bigr)
    \qquad (\alpha\in\Delta),
\]
and hence
\[
  f\bigl(p_{-\alpha}(x)\bigr)
    =p'_{-d(\alpha)}\bigl(x^{q(\alpha)}\bigr),
  \qquad
  f(w_\alpha)=w'_{d(\alpha)} .
\]
Recall (\SGARefLink{sga3:xxii:subsection:4-2}{Exposé~XXII, 4.2}) that, for every \(\alpha\in R\), every
\(z\in\mathbb G_m(S')\), and every \(t\in T(S')\),
\[
  f\bigl(\alpha^\ast(z)\bigr)
    =d(\alpha)^\ast(z)^{q(\alpha)}
    =d(\alpha)^\ast\bigl(z^{q(\alpha)}\bigr),
  \qquad
  d(\alpha)\bigl(f(t)\bigr)=\alpha(t)^{q(\alpha)} .
\]

\SGARefTarget{sga3:XXIII:1.5}
\SGAParagraph{1.5}
We call a root datum equipped with a system of simple roots a
\emph{pinned root datum}, and a \(p\)-morphism of root data
(\SGARefLink{sga3:XXI:6.8}{Exposé~XXI, 6.8}) that sends simple roots to simple roots a
\emph{\(p\)-morphism of pinned root data}.

If \(G\) is a pinned \(S\)-group, we denote by \(\mathscr R(G)\) the
corresponding pinned root datum (the root datum of \SGARefLink{sga3:xxii:prop:1.14}{Exposé~XXII, 1.14}
equipped with \(\Delta\)). Let \(p\) be the integer defined in
\SGARefLink{sga3:xxii:rem:4.2.2}{Exposé~XXII, 4.2.2}. Then:

\SGARefTarget{sga3:XXIII:1.6}
\SGAStatement{Scholium 1.6}{The correspondence
\(G\longmapsto\mathscr R(G)\) defines a functor from the category of
pinned reductive \(S\)-groups to the category of pinned root data
(whose morphisms are the \(p\)-morphisms).}

% END inputs/XXIII/tex/components/02_expose_XXIII_par14_through_scholium16.tex

% BEGIN inputs/XXIII/tex/components/03_expose_XXIII_par17_through_remark193.tex
\SGASubsection{1.7}{The pinned groups \(Z_{\Delta_1}\)}
\SGARefTarget{sga3:XXIII:1.7}

Let \(\Delta_1\) be a subset of the system \(\Delta\) of simple roots
of the pinned group \(G\). Let \(T_{\Delta_1}\) be the maximal torus
of \(\bigcap_{\alpha\in\Delta_1}\ker(\alpha)\), and set
\[
  Z_{\Delta_1}=\operatorname{Centr}_G(T_{\Delta_1}).
\]
Put \(R'=\mathbb Z\Delta_1\cap R\). It is known
(\SGARefLink{sga3:xxii:para:5.10.7}{Exposé~XXII, 5.10.7}) that \(Z_{\Delta_1}\) is a reductive \(S\)-group
with radical \(T_{\Delta_1}\), that \((T,M,R')\) is a splitting of
this group, and that \(\Delta_1\) is a system of simple roots. Thus
\[
  (Z_{\Delta_1},T,M,R',\Delta_1,(X_\alpha)_{\alpha\in\Delta_1})
\]
is a pinned \(S\)-group. We shall always give \(Z_{\Delta_1}\) this
pinning. In particular, we shall consider
\[
  Z_\alpha=Z_{\{\alpha\}},\qquad
  Z_{\alpha\beta}=Z_{\{\alpha,\beta\}}.
\]
Write \(B_{\Delta_1}=B\cap Z_{\Delta_1}\). It is known
(\textit{loc.\ cit.}) that this is the canonical Borel subgroup
\NDE{That is, the Borel subgroup of \(Z_{\Delta_1}\) corresponding to
\(R'^+=R'\cap R^+\).}
of \(Z_{\Delta_1}\), and that its unipotent part is
\(U_{\Delta_1}=U\cap Z_{\Delta_1}\). In particular,
\[
  B_\alpha=T\cdot U_\alpha.
\]
We write
\[
  U_{\alpha\beta}=U_{\{\alpha,\beta\}}
  =U\cap Z_{\alpha\beta}
  =\prod_{\gamma\in R_{\alpha\beta}^{+}}U_\gamma,
\]
where \(R_{\alpha\beta}^{+}\) is the set of positive roots that are
linear combinations of \(\alpha\) and \(\beta\).

Now let \(f:G\to G'\) be a morphism of pinned \(S\)-groups. If
\(d:R\to R'\) is the corresponding bijection and
\(\Delta_1\subset\Delta\), then
\(d(\Delta_1)=\Delta'_1\subset\Delta'\). Clearly \(f\) maps
\(T_{\Delta_1}\) into \(T'_{\Delta'_1}\), and hence
\(Z_{\Delta_1}\) into \(Z'_{\Delta'_1}\). The corresponding
\(S\)-morphism
\[
  f_{\Delta_1}:Z_{\Delta_1}\longrightarrow Z'_{\Delta'_1}
\]
is compatible with the canonical pinnings. It defines a morphism of
pinned root data
\[
  \mathscr R(f_{\Delta_1}):
  \mathscr R(Z_{\Delta_1},T,M,\ldots)
  \longrightarrow
  \mathscr R(Z'_{\Delta'_1},T',M',\ldots),
\]
and the underlying morphisms \(M'\to M\) of \(\mathscr R(f)\) and
\(\mathscr R(f_{\Delta_1})\) coincide.

\SGASubsection{1.8}{Study of the group
\(\operatorname{Norm}_G(T)\)}
\SGARefTarget{sga3:XXIII:1.8}

For every pair \((\alpha,\beta)\) of simple roots, let
\(n_{\alpha\beta}\) be the order of \(s_\alpha s_\beta\) in the Weyl
group \(W\). In particular \(n_{\alpha\alpha}=1\), and therefore
\((w_\alpha w_\beta)^{n_{\alpha\beta}}\in T(S)\).

\SGARefTarget{sga3:XXIII:1.8.1}
\SGAStatement{Definition 1.8.1}{For every
\((\alpha,\beta)\in\Delta\times\Delta\), set
\[
  t_{\alpha\beta}=(w_\alpha w_\beta)^{n_{\alpha\beta}}\in T(S).
\]
Moreover, set (\SGARefLink{sga3:xx:3.1}{Exposé~XX, 3.1})
\[
  t_\alpha=t_{\alpha\alpha}=w_\alpha^2=\alpha^\ast(-1)\in T(S).
\]}

\SGARefTarget{sga3:XXIII:1.8.2}
\SGAStatement{Proposition 1.8.2}{Let \(H\) be an \(S\)-functor in
groups that carries direct sums of schemes to products (for example,
a sheaf for the Zariski topology). Let
\[
  f_T:T\longrightarrow H
\]
be a morphism of groups, and let \(h_\alpha\in H(S)\)
(\(\alpha\in\Delta\)). There exists a necessarily unique morphism of
groups
\[
  f_N:\operatorname{Norm}_G(T)\longrightarrow H
\]
inducing \(f_T\) on \(T\) and satisfying \(f_N(w_\alpha)=h_\alpha\)
for \(\alpha\in\Delta\) if and only if the following two conditions
hold:
\begin{enumerate}
  \item[\SGAInlineRefTarget{sga3:XXIII:1.8.2:item:i}(i)] for every \(\alpha\in\Delta\),
    \[
      f_T(s_\alpha(t))=h_\alpha f_T(t)h_\alpha^{-1}
    \]
    for every \(t\in T(S')\), \(S'\to S\); equivalently,
    \(f_T\circ s_\alpha=\operatorname{int}(h_\alpha)\circ f_T\);
  \item[\SGAInlineRefTarget{sga3:XXIII:1.8.2:item:ii}(ii)] for every
    \((\alpha,\beta)\in\Delta\times\Delta\),
    \[
      f_T(t_{\alpha\beta})
        =(h_\alpha h_\beta)^{n_{\alpha\beta}}.
    \]
\end{enumerate}}

\SGAProof
Equip \((\mathrm{Sch})\) with the topology \(\mathcal T\) generated
by the pretopology whose covering families are direct sums. The
hypothesis says that \(H\) is a \(\mathcal T\)-sheaf. Let \(L\) be
the free group on the generators \((m_\alpha)_{\alpha\in R}\), and
let \(L_1\) be the invariant subgroup generated by
\[
  (m_\alpha m_\beta)^{n_{\alpha\beta}},
  \qquad(\alpha,\beta)\in\Delta\times\Delta.
\]
Let \(g:L\to W\) be defined by \(g(m_\alpha)=s_\alpha\). It is known
(\SGARefLink{sga3:XXI:5.1}{Exposé~XXI, 5.1}) that \(g\) induces an isomorphism
\(\overline g:L/L_1\xrightarrow{\sim}W\). Let \(L\) act on \(T\)
through \(g\), equivalently by \(m_\alpha\cdot t=s_\alpha(t)\).
Writing \(L_S\) for the constant group defined by \(L\), form the
semidirect product \(N=T\cdot L_S\). There is a unique morphism of
\(S\)-groups
\[
  h:N=T\cdot L_S\longrightarrow\operatorname{Norm}_G(T)
\]
such that \(h(m_\alpha)=w_\alpha\) and \(h(t)=t\) for every
\(t\in T(S')\), \(S'\to S\).

Let \(N_1\) be the invariant subgroup sheaf of \(N\) generated by
\[
  t_{\alpha\beta}^{-1}(m_\alpha m_\beta)^{n_{\alpha\beta}},
  \qquad(\alpha,\beta)\in\Delta\times\Delta.
\]
Clearly \(N_1\subset\ker h\), so \(h\) induces
\[
  h_1:N/N_1\longrightarrow\operatorname{Norm}_G(T).
\]
We prove that \(h_1\) is an isomorphism. Since \(h|_T\) is the
canonical immersion, hence a monomorphism, the canonical morphism
\(T\to N/N_1\) is also a monomorphism and induces an isomorphism
\(T\xrightarrow{\sim}TN_1/N_1\). Thus \(h_1\) induces an
isomorphism \(TN_1/N_1\xrightarrow{\sim}T\). It remains only to show
that
\[
  h_2:N/TN_1\longrightarrow\operatorname{Norm}_G(T)/T
\]
is an isomorphism. But \(TN_1\) is the invariant
\(\mathcal T\)-subgroup sheaf of \(N\) generated by \(T\) and the
\((m_\alpha m_\beta)^{n_{\alpha\beta}}\); hence it is the subgroup
sheaf generated by \(T\) and \(L_1\), namely \(T\cdot(L_1)_S\).
Therefore \(h_2\) identifies with
\[
  \overline g:L_S/(L_1)_S\longrightarrow W_S,
\]
which is an isomorphism by construction.

The proof of the proposition is now immediate. Necessity is clear.
Conversely, condition \SGARefLink{sga3:XXIII:1.8.2:item:i}{(i)} gives a morphism
\[
  u:N\longrightarrow H
\]
such that \(u(m_\alpha)=h_\alpha\) for \(\alpha\in\Delta\) and
\(u|_T=f_T\). Condition \SGARefLink{sga3:XXIII:1.8.2:item:ii}{(ii)} says that \(u\) vanishes on \(N_1\),
and the result follows.

\SGASubsection{1.9}{Faithfulness of the functor \(\mathscr R\)}
\SGARefTarget{sga3:XXIII:1.9}

\SGARefTarget{sga3:XXIII:1.9.1}
\SGAStatement{Proposition 1.9.1}{The functor \(\mathscr R\) of \SGARefLink{sga3:XXIII:1.6}{1.6} is
faithful: if
\[
  f,g:G\rightrightarrows G'
\]
are two morphisms of pinned groups such that
\(\mathscr R(f)=\mathscr R(g)\), then \(f=g\).}

Indeed, \(f\) and \(g\) agree on \(T\), on every \(U_\alpha\)
(\(\alpha\in\Delta\)), and on every \(U_{-\alpha}\)
(\(\alpha\in\Delta\)); it therefore suffices to apply:

\SGARefTarget{sga3:XXIII:1.9.2}
\SGAStatement{Lemma 1.9.2}{Let \(S\) be a scheme, let \(G\) be a
pinned reductive \(S\)-group, and let \(H\) be an \(S\)-presheaf in
groups separated for the fppf topology. Let
\[
  f,g:G\rightrightarrows H
\]
be two morphisms of \(S\)-groups. The following conditions are
equivalent:
\begin{enumerate}
  \item[\SGAInlineRefTarget{sga3:XXIII:1.9.2:item:i}(i)] \(f=g\);
  \item[\SGAInlineRefTarget{sga3:XXIII:1.9.2:item:ii}(ii)] \(f\) and \(g\) agree on \(T\), on every \(U_\alpha\)
    (\(\alpha\in\Delta\)), and on every \(U_{-\alpha}\)
    (\(\alpha\in\Delta\));
  \item[\SGAInlineRefTarget{sga3:XXIII:1.9.2:item:iii}(iii)] \(f\) and \(g\) agree on \(T\) and on every
    \(U_\alpha\) (\(\alpha\in\Delta\)), and
    \(f(w_\alpha)=g(w_\alpha)\) for every \(\alpha\in\Delta\).
\end{enumerate}}

\SGAProof
The implication \SGARefLink{sga3:XXIII:1.9.2:item:i}{(i)}\(\Rightarrow\)\SGARefLink{sga3:XXIII:1.9.2:item:ii}{(ii)} is trivial, while
\SGARefLink{sga3:XXIII:1.9.2:item:ii}{(ii)}\(\Rightarrow\)\SGARefLink{sga3:XXIII:1.9.2:item:iii}{(iii)} follows immediately from the definition of
\(w_\alpha\) in \SGARefLink{sga3:XXIII:1.2}{1.2}. It remains to prove
\SGARefLink{sga3:XXIII:1.9.2:item:iii}{(iii)}\(\Rightarrow\)\SGARefLink{sga3:XXIII:1.9.2:item:i}{(i)}. If \(\alpha\in R\), there is a sequence
\(\{\alpha_i\}\subset\Delta\) such that
\[
  \alpha=s_{\alpha_1}s_{\alpha_2}\cdots
    s_{\alpha_n}(\alpha_{n+1}),
\]
and hence
\[
  U_\alpha
    =\operatorname{int}(w_{\alpha_1}\cdots w_{\alpha_n})
       U_{\alpha_{n+1}} .
\]
Thus \(f\) and \(g\) agree on every \(U_\alpha\). They consequently
agree on \(\Omega\), and therefore agree everywhere
(\SGARefLink{sga3:xxii:prop:4.1.11}{Exposé~XXII, 4.1.11}).

\SGARefTarget{sga3:XXIII:1.9.3}
\SGAStatement{Remark 1.9.3}{If \(G\) is semisimple, the hypothesis
that \(f\) and \(g\) agree on \(T\) may be omitted in \SGARefLink{sga3:XXIII:1.9.2:item:ii}{(ii)} and \SGARefLink{sga3:XXIII:1.9.2:item:iii}{(iii)}.
Indeed, \(G\), as an fppf sheaf, is generated by the \(U_\alpha\),
\(\alpha\in R\) (\SGARefLink{sga3:xxii:rem:6.2.2:item:a}{Exposé~XXII, 6.2.2(a)}).}

% END inputs/XXIII/tex/components/03_expose_XXIII_par17_through_remark193.tex

% BEGIN inputs/XXIII/tex/components/04_expose_XXIII_theorem21_through_corollary24.tex
\SGASection{2}{Generators and relations for a pinned group}
\SGARefTarget{sga3:XXIII:2}

Throughout this section \(G\) is a fixed pinned \(S\)-group. For
\(\alpha,\beta\in\Delta\), we use the notation
\(Z_\alpha,Z_{\alpha\beta},U_{\alpha\beta},R^+_{\alpha\beta}\)
of \SGARefLink{sga3:XXIII:1.7}{1.7}.

\SGARefTarget{sga3:XXIII:2.1}
\SGAStatement{Theorem 2.1}{Let \(S\) be a scheme, let \(G\) be a
pinned \(S\)-group, and let \(H\) be an fppf \(S\)-sheaf in groups.
Suppose given morphisms of groups
\[
  f_N:\operatorname{Norm}_G(T)\longrightarrow H,\qquad
  f_\alpha:U_\alpha\longrightarrow H\quad(\alpha\in R).
\]
There exists a necessarily unique morphism of groups
\[
  f:G\longrightarrow H
\]
extending \(f_N\) and all the \(f_\alpha\), \(\alpha\in R\), if and
only if the following conditions hold:
\begin{enumerate}
  \item[\SGAInlineRefTarget{sga3:XXIII:2.1:item:i}(i)] for every \(\alpha\in\Delta\) and \(\beta\in R\),
    \[
      \operatorname{int}\bigl(f_N(w_\alpha)\bigr)\circ f_\beta
        =f_{s_\alpha(\beta)}\circ\operatorname{int}(w_\alpha);
    \]
  \item[\SGAInlineRefTarget{sga3:XXIII:2.1:item:ii}(ii)] for every \(\alpha\in\Delta\), there exists a morphism
    of groups
    \[
      F_\alpha:Z_\alpha\longrightarrow H
    \]
    extending \(f_\alpha\), \(f_{-\alpha}\), and
    \(f_N|_{\operatorname{Norm}_{Z_\alpha}(T)}\);
  \item[\SGAInlineRefTarget{sga3:XXIII:2.1:item:iii}(iii)] for every
    \((\alpha,\beta)\in\Delta\times\Delta\), there exists a morphism
    of groups \(U_{\alpha\beta}\to H\) inducing \(f_\gamma\) on
    \(U_\gamma\) for every \(\gamma\in R^+_{\alpha\beta}\)
    (so \(U_\gamma\subset U_{\alpha\beta}\)).
\end{enumerate}}

\SGASubsection{2.1.1}{Proof}
\SGARefTarget{sga3:XXIII:2.1.1}

The conditions are plainly necessary. Choose a totally ordered group
structure on \(\Gamma_0(R)\) for which the positive elements of \(R\)
are precisely \(R^+\) (\SGARefLink{sga3:XXI:3.5.6}{Exposé~XXI, 3.5.6}); every product indexed by a
subset of \(R\) will be taken in this order. Let \(f_T\) denote the
restriction of \(f_N\) to \(T\), and define a morphism of sets
\[
  f:\Omega\longrightarrow H
\]
by
\[
  f\left(
    \prod_{\alpha\in R^-}y_\alpha\cdot t\cdot
    \prod_{\alpha\in R^+}x_\alpha
  \right)
  =
  \prod_{\alpha\in R^-}f_\alpha(y_\alpha)\cdot f_T(t)\cdot
  \prod_{\alpha\in R^+}f_\alpha(x_\alpha).
\]
Every morphism satisfying the theorem must extend this map. Conversely,
any morphism of groups \(f:G\to H\) extending it also extends \(f_N\):
since it extends \(f_\alpha\) and \(f_{-\alpha}\), one has
\[
  f(w_\alpha)=F_\alpha(w_\alpha)=f_N(w_\alpha),
\]
and it extends \(f_T\) by hypothesis. By \SGARefLink{sga3:xxii:prop:4.1.11:item:ii}{Exposé~XXII, 4.1.11(ii)}, it
therefore remains to prove:

\SGARefTarget{sga3:XXIII:2.1.2}
\SGAStatement{Proposition 2.1.2}{The morphism
\(f:\Omega\to H\) defined above is \emph{generically multiplicative}:
for every \(S'\to S\) and every \(x,y\in\Omega(S')\) such that
\(xy\in\Omega(S')\), one has \(f(xy)=f(x)f(y)\).}

\SGARefTarget{sga3:XXIII:2.1.3}
\SGAStatement{Lemma 2.1.3}{Let
\(n\in\operatorname{Norm}_G(T)(S')\), \(S'\to S\), and let
\(\alpha,\beta\in R\) satisfy
\(\operatorname{int}(n)U_\alpha=U_\beta\), equivalently
\(n(\alpha)=\beta\). Then
\[
  \operatorname{int}\bigl(f_N(n)\bigr)\circ f_\alpha
    =f_\beta\circ\operatorname{int}(n).
\]}

It suffices to verify this formula for a system of generators of
\(\operatorname{Norm}_G(T)\). It holds for every \(w_\alpha\),
\(\alpha\in\Delta\), by \SGARefLink{sga3:XXIII:2.1:item:i}{2.1(i)}, so it remains to treat
\(n\in T(S')\). This is immediate from \SGARefLink{sga3:XXIII:2.1:item:ii}{2.1(ii)} when \(\alpha\) is
simple. Otherwise choose \(w\in W\) with
\(w^{-1}(\alpha)\in\Delta\). Writing \(w\) as a product of simple
reflections\NDE{``Fundamental symmetries'' was replaced by ``simple
reflections.''}, one reduces to showing that if the formula holds for
\(\alpha\) and every \(n\), it also holds for
\(w_{\alpha_0}(\alpha)\) and \(t\in T(S')\), where
\(\alpha_0\in\Delta\). By \SGARefLink{sga3:XXIII:2.1:item:i}{2.1(i)},
\[
\begin{aligned}
  \operatorname{int}\bigl(f_N(t)\bigr)\circ f_{w_{\alpha_0}(\alpha)}
   &=
  \operatorname{int}\bigl(f_N(tw_{\alpha_0})\bigr)\circ
  f_\alpha\circ\operatorname{int}(w_{\alpha_0}^{-1})\\
   &=
  f_{w_{\alpha_0}(\alpha)}\circ\operatorname{int}(tw_{\alpha_0})
    \circ\operatorname{int}(w_{\alpha_0}^{-1}).
\end{aligned}
\]

\SGARefTarget{sga3:XXIII:2.1.4}
\SGAStatement{Lemma 2.1.4}{The restriction of \(f\) to \(U\)
(respectively \(U^-\)) is a morphism of groups. In particular, this
restriction is independent of the chosen ordering of the roots.}

It suffices to prove the assertion for \(U\). By \SGARefLink{sga3:xxii:lem:5.5.8}{Exposé~XXII, 5.5.8},
it is enough to show that, for every pair \(\alpha<\beta\) of positive
roots and every \(x_\alpha\in U_\alpha(S')\),
\(x_\beta\in U_\beta(S')\), \(S'\to S\),
\[
  f_\beta(x_\beta)f_\alpha(x_\alpha)f_\beta(x_\beta^{-1})
   =f(x_\beta x_\alpha x_\beta^{-1}).
\]
By \SGARefLink{sga3:xxii:cor:5.5.2}{Exposé~XXII, 5.5.2}, there are elements
\(x_\gamma\in U_\gamma(S')\), where
\(\gamma=i\alpha+j\beta\in R\), \(i>0\), \(j\geq0\)
\NDE{The printed \(i,j>0\) was corrected to \(i>0,j\geq0\), since
\(x_\alpha\) occurs in the product on the right.}, such that
\[
  x_\beta x_\alpha x_\beta^{-1}=\prod_\gamma x_\gamma.
\]
It remains to verify
\[
  f_\beta(x_\beta)f_\alpha(x_\alpha)f_\beta(x_\beta^{-1})
   =
  \prod_{\substack{\gamma=i\alpha+j\beta\\i>0,\ j\geq0}}
  f_\gamma(x_\gamma).
\]
By \SGARefLink{sga3:XXI:3.5.4}{Exposé~XXI, 3.5.4}, there is a \(w\in W\) such that
\(w(\alpha)=\alpha_0\in\Delta\) and
\(w(\beta)\in A_{\alpha_0\beta_0}\), in the notation of \SGARefLink{sga3:XXIII:1.7}{1.7}, for
some \(\beta_0\in\Delta\). Lift \(w\) to
\(n\in\operatorname{Norm}_G(T)(S)\) by \SGARefLink{sga3:xxii:cor:3.8}{Exposé~XXII, 3.8}. It suffices
to verify the relation after conjugation by \(f_N(n)\); by \SGARefLink{sga3:XXIII:2.1.3}{2.1.3},
this reduces to \(\alpha,\beta\in R^+_{\alpha_0\beta_0}\), where
condition \SGARefLink{sga3:XXIII:2.1:item:iii}{2.1(iii)} applies.

\SGARefTarget{sga3:XXIII:2.1.5}
\SGAStatement{Lemma 2.1.5}{Let
\(n\in\operatorname{Norm}_G(T)(S)\) satisfy
\(\operatorname{int}(n)U=U^-\), so that \(n\) represents the
symmetry of the Weyl group\NDE{Relative to \(\Delta\).}
(\SGARefLink{sga3:XXI:3.6.14}{Exposé~XXI, 3.6.14}). For every \(u\in U(S')\), \(S'\to S\)
(respectively \(u\in U^-(S')\)), one has
\[
  f(nun^{-1})=f_N(n)f(u)f_N(n^{-1}).
\]}
This follows immediately from \SGARefLink{sga3:XXIII:2.1.3}{2.1.3} and \SGARefLink{sga3:XXIII:2.1.4}{2.1.4}.

\SGARefTarget{sga3:XXIII:2.1.6}
\SGAStatement{Lemma 2.1.6}{Let \(u\in B(S')\),
\(v\in B^-(S')\), and \(g\in\Omega(S')\), where \(S'\to S\). Then
\[
  f(vgu)=f(v)f(g)f(u).
\]}

Write
\[
  v=v_1t_1,\qquad g=v_2t_2u_2,\qquad u=t_3u_3,
\]
with \(v_i\in U^-(S')\), \(t_i\in T(S')\), and \(u_i\in U(S')\).
On the one hand,
\[
  f(v)f(g)f(u)
  =f(v_1)f_T(t_1)f(v_2)f_T(t_2)f(u_2)f_T(t_3)f(u_3).
\]
On the other hand,
\[
\begin{aligned}
 f(vgu)
 &=f(v_1t_1v_2t_1^{-1}\,t_1t_2t_3\,
       t_3^{-1}u_2t_3u_3)\\
 &=f(v_1\cdot t_1v_2t_1^{-1})f_T(t_1t_2t_3)
       f(t_3^{-1}u_2t_3\cdot u_3).
\end{aligned}
\]
Using \SGARefLink{sga3:XXIII:2.1.4}{2.1.4} on the two outer factors gives
\[
 f(vgu)=f(v_1)f(t_1v_2t_1^{-1})f_T(t_1t_2t_3)
   f(t_3^{-1}u_2t_3)f(u_3).
\]
It remains only to use
\[
\begin{aligned}
 f(t_1v_2t_1^{-1})&=f_T(t_1)f(v_2)f_T(t_1)^{-1},\\
 f(t_3^{-1}u_2t_3)&=f_T(t_3)^{-1}f(u_2)f_T(t_3),
\end{aligned}
\]
which follow from the definition of \(f\) and \SGARefLink{sga3:XXIII:2.1.3}{2.1.3}.

\SGARefTarget{sga3:XXIII:2.1.7}
\SGAStatement{Lemma 2.1.7}{Let \(\alpha\in\Delta\) and
\[
  g\in\Omega(S')\cap\operatorname{int}(w_\alpha)^{-1}\Omega(S'),
  \qquad S'\to S.
\]
Then
\[
  f(w_\alpha gw_\alpha^{-1})
   =f_N(w_\alpha)f(g)f_N(w_\alpha)^{-1}.
\]}

By \SGARefLink{sga3:xxii:cor:5.6.8}{Exposé~XXII, 5.6.8}, write
\[
  g=ax_{-\alpha}tx_\alpha b,
\]
where \(a\in U^-_{-\widehat\alpha}(S')\),
\(x_{-\alpha}\in U_{-\alpha}(S')\), \(t\in T(S')\),
\(x_\alpha\in U_\alpha(S')\), and
\(b\in U_{\widehat\alpha}(S')\). By \SGARefLink{sga3:XXIII:2.1.3}{2.1.3} and \SGARefLink{sga3:XXIII:2.1.4}{2.1.4}, and because
\(s_\alpha\) permutes the positive roots other than \(\alpha\)
(\SGARefLink{sga3:XXI:3.3.2}{Exposé~XXI, 3.3.2}),
\[
  \operatorname{int}(w_\alpha)a\in U^-_{-\widehat\alpha}(S'),
  \qquad
  \operatorname{int}(w_\alpha)b\in U_{\widehat\alpha}(S').
\]
The desired formula holds for \(g=a\) and \(g=b\). By \SGARefLink{sga3:XXIII:2.1.6}{2.1.6} it
therefore suffices to treat
\(g=x_{-\alpha}tx_\alpha\in Z_\alpha(S')\); in that case everything
takes place in \(Z_\alpha\), and condition \SGARefLink{sga3:XXIII:2.1:item:ii}{2.1(ii)} applies.

\SGARefTarget{sga3:XXIII:2.1.8}
\SGAStatement{Lemma 2.1.8}{Let
\(n\in\operatorname{Norm}_G(T)(S)\). For every \(S'\to S\) and every
\(g\in\Omega(S')\) such that
\(\operatorname{int}(n)g\in\Omega(S')\),
\[
  f(ngn^{-1})=f_N(n)f(g)f_N(n)^{-1}.
\]}

This is immediate for \(n\in T(S)\), by \SGARefLink{sga3:XXIII:2.1.3}{2.1.3}. The two sides define
morphisms from
\(\Omega\cap\operatorname{int}(n)^{-1}\Omega\) to \(H\). It is
enough to find an open \(V_n\subset\Omega\), containing the unit
section, such that \(\operatorname{int}(n)V_n\subset\Omega\), and on
which \(f\circ\operatorname{int}(n)\) and
\(\operatorname{int}(f_N(n))\circ f\) agree. By the structure of
\(\operatorname{Norm}_G(T)\), it suffices to show that if this is true
for \(n'\), then it is true for \(n=n'w_\alpha\),
\(\alpha\in\Delta\). Put
\[
  V_n=\Omega\cap\operatorname{int}(w_\alpha)^{-1}V_{n'}.
\]
Then
\(\operatorname{int}(n)V_n\subset\operatorname{int}(n')V_{n'}
\subset\Omega\). If \(g\in V_n(S')\), then
\[
  \operatorname{int}(n)g
   =\operatorname{int}(n')\operatorname{int}(w_\alpha)g.
\]
Now \(\operatorname{int}(w_\alpha)g\in V_{n'}(S')\), so the
induction hypothesis gives
\[
  f\bigl(\operatorname{int}(n')\operatorname{int}(w_\alpha)g\bigr)
   =\operatorname{int}(f_N(n'))f(\operatorname{int}(w_\alpha)g).
\]
\SGARefLink{sga3:XXIII:2.1.7}{Lemma 2.1.7} gives
\[
  f(\operatorname{int}(w_\alpha)g)
   =\operatorname{int}(f_N(w_\alpha))f(g),
\]
and the conclusion follows.

We now prove \SGARefLink{sga3:XXIII:2.1.2}{Proposition 2.1.2}. Let \(x,x'\in\Omega(S')\), and write
\[
  x=vtu,\qquad x'=v't'u',
\]
so \(xx'=vt(uv')t'u'\). By \SGARefLink{sga3:XXIII:2.1.6}{2.1.6} and
\(U^-(S')\Omega(S')U(S')=\Omega(S')\), it suffices to show that
\(f(uv')=f(u)f(v')\) whenever \(uv'\in\Omega(S')\). Choose
\(n\in\operatorname{Norm}_G(T)(S')\) as in \SGARefLink{sga3:XXIII:2.1.5}{2.1.5}. Then
\[
\begin{aligned}
 f(u)&=f_N(n)^{-1}f(nun^{-1})f_N(n),\\
 f(v')&=f_N(n)^{-1}f(nv'n^{-1})f_N(n),
\end{aligned}
\]
and hence
\[
 f(u)f(v')=f_N(n)^{-1}f(nun^{-1})f(nv'n^{-1})f_N(n).
\]
Since \(nun^{-1}\in U^-(S')\) and \(nv'n^{-1}\in U(S')\),
\[
 f(nun^{-1})f(nv'n^{-1})
   =f\bigl((nun^{-1})(nv'n^{-1})\bigr)
   =f(nuv'n^{-1}).
\]
Thus
\[
  f(u)f(v')=f_N(n)^{-1}f(nuv'n^{-1})f_N(n).
\]
If \(uv'\in\Omega(S')\), \SGARefLink{sga3:XXIII:2.1.8}{Lemma 2.1.8} finishes the proof.

\SGARefTarget{sga3:XXIII:2.2}
\SGAStatement{Remark 2.2}{Instead of prescribing \(f_N\), one may
prescribe a morphism \(f_T:T\to H\) and sections
\(h_\alpha\in H(S)\), \(\alpha\in\Delta\), satisfying the conditions
of \SGARefLink{sga3:XXIII:1.8.2}{1.8.2}. Condition \SGARefLink{sga3:XXIII:1.8.2:item:ii}{(ii)} of \SGARefLink{sga3:XXIII:2.1}{Theorem 2.1} must then be replaced by:
\begin{quote}
\emph{(ii\('\)) There exists a morphism of groups
\(F_\alpha:Z_\alpha\to H\) extending \(f_\alpha,f_{-\alpha},f_T\)
and satisfying \(F_\alpha(w_\alpha)=h_\alpha\).}
\end{quote}}

We now give the preceding theorem a more explicit form, retaining its
notation.

\SGARefTarget{sga3:XXIII:2.3}
\SGAStatement{Theorem 2.3}{Let \(S\) be a scheme, let \(G\) be a
pinned \(S\)-group, and let \(H\) be an fppf \(S\)-sheaf in groups.
Suppose given morphisms of groups
\[
  f_T:T\longrightarrow H,\qquad
  f_\alpha:U_\alpha\longrightarrow H\quad(\alpha\in\Delta),
\]
and sections \(h_\alpha\in H(S)\), \(\alpha\in\Delta\). There exists
a necessarily unique morphism of groups \(f:G\to H\) inducing
\(f_T\) and the \(f_\alpha\) and satisfying
\(f(w_\alpha)=h_\alpha\) for all \(\alpha\in\Delta\) if and only if:
\begin{enumerate}
  \item[\SGAInlineRefTarget{sga3:XXIII:2.3:item:i}(i)] for every \(S'\to S\), \(t\in T(S')\),
    \(\alpha\in\Delta\), and \(x\in U_\alpha(S')\),
    \begin{equation}
    \tag{1}\label{sga3:XXIII:2.3:eq1}
      f_T(t)f_\alpha(x)f_T(t)^{-1}
       =f_\alpha(\operatorname{int}(t)x);
    \end{equation}
  \item[\SGAInlineRefTarget{sga3:XXIII:2.3:item:ii}(ii)] for every \(\alpha\in\Delta\), \(S'\to S\), and
    \(t\in T(S')\),
    \begin{equation}
    \tag{2}\label{sga3:XXIII:2.3:eq2}
      h_\alpha f_T(t)h_\alpha^{-1}
       =f_T(s_\alpha(t))
       =f_T\!\bigl(t\alpha^\ast(\alpha(t)^{-1})\bigr);
    \end{equation}
  \item[(iii)] for every
    \((\alpha,\beta)\in\Delta\times\Delta\),
    \begin{equation}
    \tag{3}\label{sga3:XXIII:2.3:eq3}
      (h_\alpha h_\beta)^{n_{\alpha\beta}}
       =f_T(t_{\alpha\beta});
    \end{equation}
  \item[(iv)] for every \(\alpha\in\Delta\), recalling that
    \(u_\alpha=\exp_\alpha(X_\alpha)\),
    \begin{equation}
    \tag{4}\label{sga3:XXIII:2.3:eq4}
      (h_\alpha f_\alpha(u_\alpha))^3=e;
    \end{equation}
  \item[\SGAInlineRefTarget{sga3:XXIII:2.3:item:v}(v)] for every distinct
    \(\alpha,\beta\in\Delta\), there exists a morphism
    \(f_{\alpha\beta}:U_{\alpha\beta}\to H\) such that
    \begin{equation}
    \tag{5}\label{sga3:XXIII:2.3:eq5}
      f_{\alpha\beta}|_{U_\alpha}=f_\alpha,\qquad
      f_{\alpha\beta}|_{U_\beta}=f_\beta,
    \end{equation}
    and, for every \(\gamma\in R^+_{\alpha\beta}\) with
    \(\gamma\ne\alpha\) (respectively \(\gamma\ne\beta\)) and every
    \(x\in U_\gamma(S')\), \(S'\to S\),
    \begin{equation}
    \tag{6}\label{sga3:XXIII:2.3:eq6}
      \operatorname{int}(h_\alpha)f_{\alpha\beta}(x)
       =f_{\alpha\beta}(\operatorname{int}(w_\alpha)x)
    \end{equation}
    (respectively
    \(\operatorname{int}(h_\beta)f_{\alpha\beta}(x)
     =f_{\alpha\beta}(\operatorname{int}(w_\beta)x)\)).
\end{enumerate}}

\SGAProof
Uniqueness follows from \SGARefLink{sga3:XXIII:1.9.2}{1.9.2}. We prove existence.

\SGARefTarget{sga3:XXIII:2.3.1}
\SGAStatement{Lemma 2.3.1}{There exists a morphism
\[
  f_N:\operatorname{Norm}_G(T)\longrightarrow H
\]
extending \(f_T\) and satisfying \(f_N(w_\alpha)=h_\alpha\).}
This is precisely \SGARefLink{sga3:XXIII:1.8.2}{1.8.2}, in view of conditions
\eqref{sga3:XXIII:2.3:eq2} and \eqref{sga3:XXIII:2.3:eq3}.

\SGARefTarget{sga3:XXIII:2.3.2}
\SGAStatement{Lemma 2.3.2}{For every \(\alpha\in\Delta\), there is a
morphism of groups \(F_\alpha:Z_\alpha\to H\) extending \(f_T\) and
\(f_\alpha\), satisfying \(F_\alpha(w_\alpha)=h_\alpha\), and hence
extending \(f_N|_{\operatorname{Norm}_{Z_\alpha}(T)}\).}
This follows from \SGARefLink{sga3:xx:6.2}{Exposé~XX, 6.2} and conditions
\eqref{sga3:XXIII:2.3:eq1}, \eqref{sga3:XXIII:2.3:eq3}, and
\eqref{sga3:XXIII:2.3:eq4}.

\SGARefTarget{sga3:XXIII:2.3.3}
\SGAStatement{Lemma 2.3.3}{Let
\((\alpha,\beta)\in\Delta\times\Delta\), \(\alpha\ne\beta\). Let
\(n\in\operatorname{Norm}_{Z_{\alpha\beta}}(T)(S)\) satisfy
\(n(\alpha)=\alpha\) (respectively \(n(\alpha)=\beta\)), equivalently
\(\operatorname{int}(n)U_\alpha=U_\alpha\) (respectively
\(U_\beta\)). Then, for every \(S'\to S\) and \(x\in U_\alpha(S')\),
\[
  \operatorname{int}(f_N(n))f_\alpha(x)
   =f_\alpha(\operatorname{int}(n)x)
\]
(respectively the right-hand side is
\(f_\beta(\operatorname{int}(n)x)\)).}

There exist \(t\in T(S)\) and a sequence
\(\{\alpha_i\}\subset\{\alpha,\beta\}\) such that
\[
  n=tw_{\alpha_k}\cdots w_{\alpha_1}.
\]
Condition \eqref{sga3:XXIII:2.3:eq1} treats \(n=t\), so we may omit
\(t\). Proceed by induction on \(k\), and set
\[
  \gamma_i=s_{\alpha_i}\cdots s_{\alpha_1}(\alpha).
\]
If every \(\gamma_i\) is positive, condition \SGARefLink{sga3:XXIII:2.3:item:v}{(v)} yields
\[
  \operatorname{int}\bigl(f_N(w_{\alpha_i}\cdots w_{\alpha_1})\bigr)
     f_\alpha(x)
   =
  f_{\alpha\beta}\bigl(
     \operatorname{int}(w_{\alpha_i}\cdots w_{\alpha_1})x\bigr),
\]
and \(i=k\) gives the result. Otherwise there is a \(j<k\) such that
\(\gamma_j\in R^+\) but \(\gamma_{j+1}\notin R^+\). Since
\(\gamma_{j+1}=s_{\alpha_j}(\gamma_j)\), \SGARefLink{sga3:XXI:3.3.2}{Exposé~XXI, 3.3.2} gives
\(\gamma_j=\alpha_j\), hence \(\gamma_j\) is \(\alpha\) or \(\beta\).
Decompose
\[
 n=n'n'',\qquad
 n'=w_{\alpha_k}\cdots w_{\alpha_{j+1}},\qquad
 n''=w_{\alpha_j}\cdots w_{\alpha_1}.
\]
Both \(n'\) and \(n''\) satisfy the lemma's hypotheses and involve
fewer than \(k\) reflections, so induction applies.

\SGARefTarget{sga3:XXIII:2.3.4}
\SGAStatement{Lemma 2.3.4}{Let \(\alpha\in R\). Suppose
\(n,n'\in\operatorname{Norm}_G(T)(S)\) and
\(\beta,\beta'\in\Delta\) satisfy
\(n(\alpha)=\beta\) and \(n'(\alpha)=\beta'\). Then, for every
\(x\in U_\alpha(S')\), \(S'\to S\),
\[
 \operatorname{int}(f_N(n)^{-1})f_\beta(nx n^{-1})
 =
 \operatorname{int}(f_N(n')^{-1})f_{\beta'}(n'x n'^{-1}).
\]}

It suffices to prove that if \(n(\alpha)=\beta\), with
\(\alpha,\beta\in\Delta\), then
\[
  \operatorname{int}(f_N(n))\circ f_\alpha
    =f_\beta\circ\operatorname{int}(n).
\]
By Tits's lemma (\SGARefLink{sga3:XXI:5.6}{Exposé~XXI, 5.6}), there are simple roots
\[
  \alpha_0=\alpha,\alpha_1,\ldots,\alpha_m=\beta
\]
and elements \(w_i\in W\), \(0\leq i<m\), such that
\[
  n=w_{m-1}\cdots w_0,\qquad
  w_i(\alpha_i)=\alpha_{i+1},
\]
and for each \(i\) there is a simple root \(\beta_i\) such that
\[
  w_i\in W_{\alpha_i\beta_i},\qquad
  \alpha_{i+1}=\alpha_i\ \text{or}\ \beta_i.
\]
Induction reduces the assertion to \SGARefLink{sga3:XXIII:2.3.3}{Lemma 2.3.3}.

\SGARefTarget{sga3:XXIII:2.3.5}
\SGAStatement{Lemma 2.3.5}{There is a family of morphisms of groups
\(f'_\gamma:U_\gamma\to H\), \(\gamma\in R\), such that:
\begin{enumerate}
  \item[\SGAInlineRefTarget{sga3:XXIII:2.3.5:item:i}(i)] for \(\alpha\in\Delta\),
    \(f'_\alpha=f_\alpha\) and
    \(f'_{-\alpha}=F_\alpha|_{U_{-\alpha}}\);
  \item[\SGAInlineRefTarget{sga3:XXIII:2.3.5:item:ii}(ii)] if \(\alpha,\beta\in\Delta\) and
    \(\gamma\in R^+_{\alpha\beta}\), then
    \(f'_\gamma=f_{\alpha\beta}|_{U_\gamma}\);
  \item[\SGAInlineRefTarget{sga3:XXIII:2.3.5:item:iii}(iii)] if \(n\in\operatorname{Norm}_G(T)(S)\) and
    \(\alpha,\beta\in R\) satisfy \(n(\alpha)=\beta\), then
    \[
      \operatorname{int}(f_N(n))f'_\alpha(x)
       =f'_\beta(\operatorname{int}(n)x)
    \]
    for every \(x\in U_\alpha(S')\), \(S'\to S\).
\end{enumerate}}

For every root \(\alpha\), choose
\(n\in\operatorname{Norm}_G(T)(S)\) with \(n(\alpha)\in\Delta\), and
define \(f'_\alpha(x)\) by the expression in \SGARefLink{sga3:XXIII:2.3.4}{Lemma 2.3.4}. This is
independent of \(n\), and \(f'_\alpha\) is a morphism of groups.
Property \SGARefLink{sga3:XXIII:2.3.5:item:iii}{(iii)} holds by construction. The first part of \SGARefLink{sga3:XXIII:2.3.5:item:i}{(i)} follows
by taking \(n=e\), and the second by taking \(n=w_\alpha\) and using
\SGARefLink{sga3:XXIII:2.3.2}{2.3.2}. If \(\gamma\in R^+_{\alpha\beta}\), choose
\(n\in\operatorname{Norm}_{Z_{\alpha\beta}}(T)(S)\) with
\(n(\gamma)=\alpha\) or \(\beta\); then \SGARefLink{sga3:XXIII:2.3.5:item:iii}{(iii)}, together with
conditions \eqref{sga3:XXIII:2.3:eq5} and
\eqref{sga3:XXIII:2.3:eq6}, proves \SGARefLink{sga3:XXIII:2.3.5:item:ii}{(ii)}.

\SGARefTarget{sga3:XXIII:2.3.6}
\SGAParagraph{2.3.6}
We finish the proof by verifying the conditions of \SGARefLink{sga3:XXIII:2.1}{Theorem 2.1} for
\(f_N\) and \(f'_\alpha\), \(\alpha\in R\):
\begin{itemize}
  \item \SGARefLink{sga3:XXIII:2.1:item:i}{2.1(i)} follows from \SGARefLink{sga3:XXIII:2.3.5:item:iii}{2.3.5(iii)};
  \item \SGARefLink{sga3:XXIII:2.1:item:ii}{2.1(ii)} follows from \SGARefLink{sga3:XXIII:2.3.5:item:i}{2.3.5(i)} and \SGARefLink{sga3:XXIII:2.3.2}{2.3.2};
  \item \SGARefLink{sga3:XXIII:2.1:item:iii}{2.1(iii)} follows from \SGARefLink{sga3:XXIII:2.3.5:item:ii}{2.3.5(ii)} and the fact that
    \(f_{\alpha\beta}\) is a morphism of groups.
\end{itemize}

An immediate corollary is:

\SGARefTarget{sga3:XXIII:2.4}
\SGAStatement{Corollary 2.4}{Let \(S\) be a scheme, let \(G\) be a
pinned \(S\)-group of semisimple rank at least \(1\), and let \(H\)
be an fppf \(S\)-sheaf in groups. For each
\((\alpha,\beta)\in\Delta\times\Delta\), let
\[
  F_{\alpha\beta}:Z_{\alpha\beta}\longrightarrow H
\]
be a morphism of groups. There exists a morphism \(f:G\to H\)
inducing the \(F_{\alpha\beta}\) if and only if, for every
\((\alpha,\beta)\in\Delta\times\Delta\),
\[
  F_{\alpha\beta}|_{Z_\alpha}=F_{\alpha\alpha}.
\]}

Necessity is clear. Conversely, put
\[
  f_T=F_{\alpha\alpha}|_T;
\]
this is independent of \(\alpha\), since
\[
  F_{\alpha\alpha}|_T
   =F_{\alpha\beta}|_T
   =F_{\beta\beta}|_T.
\]
Set
\[
  f_\alpha=F_{\alpha\alpha}|_{U_\alpha},\qquad
  h_\alpha=F_{\alpha\alpha}(w_\alpha),\qquad
  f_{\alpha\beta}=F_{\alpha\beta}|_{U_{\alpha\beta}}.
\]
The conditions of \SGARefLink{sga3:XXIII:2.3}{Theorem 2.3} hold: (1), (2), and (4) hold in
\(Z_\alpha\), while (3), (5), and (6) hold in \(Z_{\alpha\beta}\).
Thus there is a morphism \(f:G\to H\) extending \(f_T\) and all
\(f_\alpha\), and satisfying \(f(w_\alpha)=h_\alpha\). It agrees
with \(F_{\alpha\beta}\) on the generators of \(Z_{\alpha\beta}\);
hence \(f|_{Z_{\alpha\beta}}=F_{\alpha\beta}\).

% END inputs/XXIII/tex/components/04_expose_XXIII_theorem21_through_corollary24.tex

% BEGIN inputs/XXIII/tex/components/05_expose_XXIII_cor25_through_par313.tex
\SGARefTarget{sga3:XXIII:2.5}
\SGAStatement{Corollary 2.5}{Let \(S\) be a scheme; let \(G\) and
\(G'\) be two pinned \(S\)-groups of semisimple rank \(2\); let
\(q>0\) be an integer such that \(x\mapsto x^q\) is an endomorphism
of \(\mathbf G_{a,S}\); and let
\[
  h:R(G')\longrightarrow R(G)
\]
be a \(q\)-morphism of pinned root data. For each
\(\alpha\in R^+\), choose
\[
  X_\alpha\in\Gamma(S,\mathfrak g^\alpha)^\times,
  \qquad
  X'_{d(\alpha)}
    \in\Gamma(S,\mathfrak g'^{\,d(\alpha)})^\times,
\]
extending the canonical choices for \(\alpha\in\Delta\). Suppose:
\begin{enumerate}
  \item[\SGAInlineRefTarget{sga3:XXIII:2.5:item:i}(i)] if \(\Delta=\{\alpha,\beta\}\), then
    \[
      D_S(h)(t_{\alpha\beta})=t'_{d(\alpha)d(\beta)};
    \]
  \item[\SGAInlineRefTarget{sga3:XXIII:2.5:item:ii}(ii)] for every \(\alpha\in\Delta\) and every
    \(\beta\in R^+\), \(\beta\ne\alpha\), if
    \(z\in\mathbf G_m(S)\) is defined by
    \[
      \operatorname{Ad}(w_\alpha)X_\beta
        =zX_{s_\alpha(\beta)},
    \]
    then
    \[
      \operatorname{Ad}(w'_{d(\alpha)})X'_{d(\beta)}
        =z^{q(\beta)}X'_{d(s_\alpha(\beta))};
    \]
  \item[(iii)] there is a morphism of groups \(f:U\to U'\) such
    that, for every \(\alpha\in R^+\), every \(S'\to S\), and every
    \(x\in\mathbf G_a(S')\),
    \[
      f(\exp(xX_\alpha))
        =\exp\!\bigl(x^{q(\alpha)}X'_{d(\alpha)}\bigr).
    \]
\end{enumerate}
Then there is a morphism of pinned groups \(g:G\to G'\) such that
\(R(g)=h\).}

Indeed, define \(f_\alpha:U_\alpha\to G'\) by
\[
  f_\alpha(\exp(xX_\alpha))
    =\exp\!\bigl(x^{q(\alpha)}X'_{d(\alpha)}\bigr),
\]
and put
\[
  f_T=D_S(h),\qquad h_\alpha=w'_{d(\alpha)}.
\]
The conditions of \SGARefLink{sga3:XXIII:2.3}{Theorem 2.3} hold: note that
\[
  q(s_\alpha(\beta))=q(\beta)
\]
by \SGARefLink{sga3:XXI:6.8.4}{Exposé XXI, 6.8.4}, and that
\[
  D_S(h)(t_\alpha)=t'_{d(\alpha)}
\]
always holds. The assertion follows at once.

\SGARefTarget{sga3:XXIII:2.6}
\SGAStatement{Remark 2.6}{Condition (v) of \SGARefLink{sga3:XXIII:2.3}{Theorem 2.3} may be made
more explicit as follows. One must first verify:
\begin{enumerate}
  \item[\SGAInlineRefTarget{sga3:XXIII:2.6:item:a}(a)] for every word in \(w_\alpha,w_\beta\) whose
    corresponding word carries \(\alpha\) or \(\beta\) to
    \(\alpha\) or \(\beta\), the corresponding relation of the type
    occurring in \SGARefLink{sga3:XXIII:2.3.3}{2.3.3} holds.
\end{enumerate}
The proof of \SGARefLink{sga3:XXIII:2.3.3}{2.3.3} shows that it suffices to check the words that are
minimal, in the sense that no nontrivial initial subword satisfies
the imposed conditions.

If \SGARefLink{sga3:XXIII:2.6:item:a}{(a)} holds, define \(f_\gamma:U_\gamma\to H\) for every
\(\gamma\in R^+_{\alpha\beta}\) as in \SGARefLink{sga3:XXIII:2.3.5}{2.3.5}. One must then require:
\begin{enumerate}
  \item[\SGAInlineRefTarget{sga3:XXIII:2.6:item:b}(b)] the map defined by the \(f_\gamma\),
    \[
      U_{\alpha\beta}
        =\prod_{\gamma\in R^+_{\alpha\beta}}U_\gamma
        \longrightarrow H,
    \]
    is a morphism of groups.
\end{enumerate}
By \SGARefLink{sga3:xxii:lem:5.5.8}{Exposé XXII, 5.5.8}, condition \SGARefLink{sga3:XXIII:2.6:item:b}{(b)} follows from:
\begin{enumerate}
  \item[(b\('\))] the preceding map respects the commutation
    relations between \(U_\gamma\) and \(U_\delta\), for
    \(\gamma,\delta\in R^+_{\alpha\beta}\), \(\gamma>\delta\);
    namely, the relations involving
    \(C_{i,j,\gamma,\delta}\) in \SGARefLink{sga3:xxii:cor:5.5.2}{Exposé XXII, 5.5.2}.
\end{enumerate}}

Conversely, the conjunction of \SGARefLink{sga3:XXIII:2.6:item:a}{(a)} and (b\('\)) is equivalent to
condition (v).

These conditions may even be reduced to conditions involving only
\[
  h_\alpha,\quad h_\beta,\quad
  f_\alpha(u_\alpha),\quad f_\beta(u_\beta)
\]
in \(H\). For example, if
\(s_\alpha s_\beta s_\alpha(\beta)=\alpha\), a condition of type
(a) reads
\begin{equation}
  \operatorname{int}(h_\alpha h_\beta h_\alpha)f_\beta(x)
    =f_\alpha\!\left(
       \operatorname{int}(w_\alpha w_\beta w_\alpha)x
     \right).
  \tag{1}
  \label{sga3:XXIII:2.6:eq1}
\end{equation}
Here \(x\in U_\beta(S')\), \(S'\to S\). In particular, for
\(x=u_\beta\),
\[
  \operatorname{int}(w_\alpha w_\beta w_\alpha)u_\beta
    =u_\alpha^z
\]
for a certain \(z\in\mathbf G_m(S)\), and
\begin{equation}
  \operatorname{int}(h_\alpha h_\beta h_\alpha)f_\beta(u_\beta)
    =f_\alpha(u_\alpha)^z.
  \tag{1\('\)}
  \label{sga3:XXIII:2.6:eq1prime}
\end{equation}

Conversely, taking conditions \SGARefLink{sga3:XXIII:2.5:item:i}{(i)}--(iv) of \SGARefLink{sga3:XXIII:2.3}{2.3} into account,
\eqref{sga3:XXIII:2.6:eq1prime} implies
\eqref{sga3:XXIII:2.6:eq1}. If \(t\in T(S')\), conjugating
\eqref{sga3:XXIII:2.6:eq1prime} by \(t\), and using \SGARefLink{sga3:XXIII:2.3:item:i}{(i)} and \SGARefLink{sga3:XXIII:2.3:item:ii}{(ii)},
gives \eqref{sga3:XXIII:2.6:eq1} for
\[
  x=\operatorname{int}(t)u_\beta=u_\beta^{\beta(t)}.
\]
The morphism \(\beta:T\to\mathbf G_{m,S}\) is faithfully flat, so
\eqref{sga3:XXIII:2.6:eq1} holds for
\(x\in U_\beta(S')^\times\). It is additive in \(x\), and every
section of \(U_\beta\) is locally the sum of two invertible sections.
Thus
\[
  \text{(1')}+\text{(i)}+\text{(ii)}
    \quad\Longrightarrow\quad\text{(1)}.
\]

The same argument applies to conditions of type \SGARefLink{sga3:XXIII:2.6:item:b}{(b)}. One uses the
fact that, for two distinct positive roots \(\gamma,\gamma'\), hence
linearly independent over \(\mathbf Q\), the morphism
\[
  T\longrightarrow\mathbf G_{m,S}^2
\]
with components \(\gamma,\gamma'\) is faithfully flat. We leave the
details of this transposition to the reader.

\SGANextSectionTarget{sga3:xxiii:section:3}\SGASection{3}{Groups of Semisimple Rank \(2\)}
\SGANextSectionTarget{sga3:xxiii:subsection:3-1}\SGASubsection{3.1}{Generalities}

\SGARefTarget{sga3:XXIII:3.1.1}
\SGAStatement{Lemma 3.1.1}{Let \(S\) be a scheme, let \(G\) be a
pinned \(S\)-group, and let \(\alpha,\beta\) be two roots of \(G\)
such that \(\alpha+\beta\ne0\).
\begin{enumerate}
  \item[\SGAInlineRefTarget{sga3:xxiii:subsection:3-1:item:i}\SGAInlineRefTarget{sga3:xxiii:lemma:3-1-1:item:i}(i)] If \(\alpha+\beta\notin R\), then
    \[
      \exp(X_\alpha)\exp(X_\beta)
        =\exp(X_\beta)\exp(X_\alpha)
    \]
    for every \(S'\to S\),
    \(X_\alpha\in W(\mathfrak g^\alpha)(S')\), and
    \(X_\beta\in W(\mathfrak g^\beta)(S')\).
  \item[\SGAInlineRefTarget{sga3:xxiii:lemma:3-1-1:item:ii}\SGAInlineRefTarget{sga3:xxiii:subsection:3-1:item:ii}(ii)] If neither \(\alpha+\beta\) nor \(\alpha-\beta\) is a
    root, then
    \[
      w_\alpha(z_\alpha)\exp(X_\beta)w_\alpha(z_\alpha)^{-1}
        =\exp(X_\beta)
    \]
    for every \(X_\beta\in W(\mathfrak g^\beta)(S')\) and
    \(z_\alpha\in W(\mathfrak g^\alpha)^\times(S')\), and
    \[
      w_\alpha(z_\alpha)w_\beta(z_\beta)
        =w_\beta(z_\beta)w_\alpha(z_\alpha)
    \]
    for every
    \(z_\alpha\in W(\mathfrak g^\alpha)^\times(S')\) and
    \(z_\beta\in W(\mathfrak g^\beta)^\times(S')\).
  \item[\SGAInlineRefTarget{sga3:XXIII:3.1.1:item:iii}\SGAInlineRefTarget{sga3:xxiii:lemma:3-1-1:item:iii}(iii)] Let
    \(X_\alpha\in W(\mathfrak g^\alpha)^\times(S')\),
    \(X_\beta\in W(\mathfrak g^\beta)^\times(S')\), and
    \(w\in\operatorname{Norm}_G(T)(S')\) satisfy
    \(w(\alpha)=\beta\). Define \(z\in\mathbf G_m(S')\) by
    \[
      \operatorname{Ad}(w)X_\alpha=zX_\beta.
    \]
    Then
    \[
    \begin{aligned}
      \operatorname{int}(w)\exp(xX_\alpha)
        &=\exp(xzX_\beta),\\
      \operatorname{int}(w)\exp(yX_\alpha^{-1})
        &=\exp(yz^{-1}X_\beta^{-1}),\\
      \operatorname{int}(w)w_\alpha(X_\alpha)
        &=\beta^*(z)w_\beta(X_\beta).
    \end{aligned}
    \]
    In particular, if \(z=\pm1\), then
    \[
      \operatorname{int}(w)w_\alpha(X_\alpha)
        =w_\beta(X_\beta)^z.
    \]
  \item[\SGAInlineRefTarget{sga3:XXIII:3.1.1:item:iv}(iv)] If
    \(t_\alpha=\alpha^*(-1)\) and \(t_\beta=\beta^*(-1)\), then
    \[
      s_\alpha(t_\beta)
        =t_\beta t_\alpha^{(\beta^*,\alpha)},
      \qquad
      \beta(t_\alpha)=(-1)^{(\alpha^*,\beta)}.
    \]
\end{enumerate}}

\SGAProof
Part \SGARefLink{sga3:xxiii:subsection:3-1:item:i}{(i)} follows immediately from \SGARefLink{sga3:xxii:cor:5.5.2}{Exposé XXII, 5.5.2}. Part \SGARefLink{sga3:xxiii:subsection:3-1:item:ii}{(ii)}
follows from \SGARefLink{sga3:xx:3.1}{Exposé XX, 3.1} and from \SGARefLink{sga3:xxiii:subsection:3-1:item:i}{(i)} applied to
\[
  (\beta,\alpha),\quad(\beta,-\alpha),\quad
  (-\beta,\alpha),\quad(-\beta,-\alpha).
\]
Part \SGARefLink{sga3:XXIII:3.1.1:item:iii}{(iii)} is immediate from the definitions; for its last assertion,
observe that
\[
  \beta^*(-1)=w_\beta(X_\beta)^{-2}.
\]
Finally, \SGARefLink{sga3:XXIII:3.1.1:item:iv}{(iv)} is tautological.

\SGARefTarget{sga3:XXIII:3.1.2}
\SGAStatement{Proposition 3.1.2 (Groups of type
\(A_1\times A_1\))}{Let \(S\) be a scheme and \(G\) a pinned
\(S\)-group of type \(A_1\times A_1\). Write
\(\Delta=R^+=\{\alpha,\beta\}\).
\begin{enumerate}
  \item[\SGAInlineRefTarget{sga3:XXIII:3.1.2:item:i}(i)]
    \[
      t_{\alpha\beta}=(w_\alpha w_\beta)^2
        =t_\alpha t_\beta
        =(w_\beta w_\alpha)^2=t_{\beta\alpha}.
    \]
  \item[\SGAInlineRefTarget{sga3:XXIII:3.1.2:item:ii}(ii)]\SGAInlineRefTarget{sga3:xxiii:proposition:3-1-2:item:ii}
    \[
      \operatorname{Ad}(w_\alpha)X_\beta=X_\beta,
      \qquad
      \operatorname{Ad}(w_\beta)X_\alpha=X_\alpha.
    \]
  \item[\SGAInlineRefTarget{sga3:XXIII:3.1.2:item:iii}(iii)] The groups \(U_\alpha\) and \(U_\beta\) commute; in
    other words, \(U\) is commutative.
\end{enumerate}}

\SGAProof
\SGARefLink{sga3:xxiii:lemma:3-1-1:item:ii}{Lemma 3.1.1(ii)} gives \(w_\alpha w_\beta=w_\beta w_\alpha\), whence
\[
  (w_\alpha w_\beta)^2=w_\alpha^2w_\beta^2=t_\alpha t_\beta,
\]
which proves \SGARefLink{sga3:XXIII:3.1.2:item:i}{(i)}. The same lemma gives \SGARefLink{sga3:XXIII:3.1.2:item:ii}{(ii)}, and \SGARefLink{sga3:XXIII:3.1.2:item:iii}{(iii)} is
\SGARefLink{sga3:xxiii:lemma:3-1-1:item:i}{Lemma 3.1.1(i)}.

\SGARefTarget{sga3:XXIII:3.1.3}
\SGAParagraph{3.1.3}
Condition (v) of \SGARefLink{sga3:XXIII:2.3}{Theorem 2.3} takes the following explicit form. Put
\[
  v_\alpha=f_\alpha(u_\alpha)
  \qquad(\alpha\in\Delta).
\]
Then the two groups of conditions are
\[
\text{(A)}\quad
\begin{cases}
  h_\alpha v_\beta h_\alpha^{-1}=v_\beta,\\
  h_\beta v_\alpha h_\beta^{-1}=v_\alpha,
\end{cases}
\qquad
\text{(B)}\quad
v_\alpha v_\beta=v_\beta v_\alpha.
\]

% END inputs/XXIII/tex/components/05_expose_XXIII_cor25_through_par313.tex

% BEGIN inputs/XXIII/tex/components/06_expose_XXIII_type_A2_and_B2.tex
\SGANextSectionTarget{sga3:xxiii:subsection:3-2}\SGASubsection{3.2}{Groups of Type \(A_2\)}

\SGARefTarget{sga3:XXIII:3.2.1}
\SGAStatement{Proposition 3.2.1}{Let \(S\) be a scheme and \(G\) a
pinned \(S\)-group of type \(A_2\). Write
\[
  \Delta=\{\alpha,\beta\},
  \qquad
  R^+=\{\alpha,\beta,\alpha+\beta\}.
\]
\begin{enumerate}
  \item[\SGAInlineRefTarget{sga3:XXIII:3.2.1:item:i}\SGAInlineRefTarget{sga3:xxiii:proposition:3-2-1:item:i}(i)]
    \[
      t_{\alpha\beta}=(w_\alpha w_\beta)^3
        =e=(w_\beta w_\alpha)^3=t_{\beta\alpha}.
    \]
  \item[\SGAInlineRefTarget{sga3:xxiii:proposition:3-2-1:item:ii}\SGAInlineRefTarget{sga3:XXIII:3.2.1:item:ii}(ii)] Put
    \(\operatorname{Ad}(w_\beta)X_\alpha=X_{\alpha+\beta}\). Then
    \[
    \begin{aligned}
      \operatorname{Ad}(w_\alpha)X_\beta&=-X_{\alpha+\beta},\\
      \operatorname{Ad}(w_\alpha)X_{\alpha+\beta}&=X_\beta,\\
      \operatorname{Ad}(w_\beta)X_{\alpha+\beta}&=-X_\alpha.
    \end{aligned}
    \]
  \item[\SGAInlineRefTarget{sga3:XXIII:3.2.1:item:iii}(iii)] Put
    \[
      p_{\alpha+\beta}(x)
        =\exp(xX_{\alpha+\beta})
        =\operatorname{int}(w_\beta)p_\alpha(x).
    \]
    Then
    \[
      p_\beta(y)p_\alpha(x)
        =p_\alpha(x)p_\beta(y)p_{\alpha+\beta}(xy).
    \]
\end{enumerate}}

\SGARefTarget{sga3:XXIII:3.2.2}
\SGAParagraph{3.2.2}
The proof occupies \SGARefLink{sga3:XXIII:3.2.2}{3.2.2}--\SGARefLink{sga3:XXIII:3.2.7}{3.2.7}. First,
\[
  (\beta^*,\alpha)=(\alpha^*,\beta)=-1,
\]
so \(\alpha(t_\beta)=\beta(t_\alpha)=-1\). Put
\[
  \operatorname{Ad}(w_\beta)X_\alpha=X_{\alpha+\beta}.
\]
Then
\[
  \operatorname{Ad}(w_\beta)X_{\alpha+\beta}
    =\alpha(t_\beta)X_\alpha=-X_\alpha.
\]
Write
\[
  \operatorname{Ad}(w_\alpha)X_\beta=zX_{\alpha+\beta},
  \qquad z\in\mathbf G_m(S).
\]
It follows that
\[
  \operatorname{Ad}(w_\alpha)X_{\alpha+\beta}
    =\beta(t_\alpha)z^{-1}X_\beta=-z^{-1}X_\beta.
\]
By \SGARefLink{sga3:xxii:cor:5.5.2}{Exposé XXII, 5.5.2}, there is a unique
\(A\in\mathbf G_a(S)\) such that
\begin{equation}
  p_\beta(y)p_\alpha(x)
    =p_\alpha(x)p_\beta(y)p_{\alpha+\beta}(Axy).
  \tag{+}
  \label{sga3:XXIII:3.2:eqplus}
\end{equation}
Thus \SGARefLink{sga3:XXIII:3.2.1:item:ii}{(ii)} and \SGARefLink{sga3:XXIII:3.2.1:item:iii}{(iii)} reduce to proving \(A=-z=1\).

\SGARefTarget{sga3:XXIII:3.2.3}
\SGAParagraph{3.2.3}
Conjugating \eqref{sga3:XXIII:3.2:eqplus} by \(w_\beta\) gives
\begin{equation}
  p_{-\beta}(-y)p_{\alpha+\beta}(x)
    =p_{\alpha+\beta}(x)p_{-\beta}(-y)p_\alpha(-Axy).
  \tag{++}
  \label{sga3:XXIII:3.2:eqplusplus}
\end{equation}

\SGARefTarget{sga3:XXIII:3.2.4}
\SGAParagraph{3.2.4}
By definition,
\[
  w_\beta p_\alpha(x)w_\beta^{-1}=p_{\alpha+\beta}(x).
\]
Expanding \(w_\beta\), this is
\[
\begin{aligned}
  &p_\beta(1)p_{-\beta}(-1)p_\beta(1)p_\alpha(x)
    p_\beta(-1)p_{-\beta}(1)p_\beta(-1)\\
  &\hspace{55mm}=p_{\alpha+\beta}(x).
\end{aligned}
\]
Since \(p_\beta\) commutes with \(p_{\alpha+\beta}\), this may be
written
\[
  p_\beta(1)p_\alpha(x)p_\beta(-1)
    =p_{-\beta}(1)p_{\alpha+\beta}(x)p_{-\beta}(-1).
\]
Using \eqref{sga3:XXIII:3.2:eqplus} on the left and
\eqref{sga3:XXIII:3.2:eqplusplus} on the right gives
\[
\begin{aligned}
  p_\alpha(x)p_\beta(1)p_{\alpha+\beta}(Ax)p_\beta(-1)
   ={}&p_{\alpha+\beta}(x)p_{-\beta}(1)\\
     &{}\cdot p_\alpha(Ax)p_{-\beta}(-1).
\end{aligned}
\]
Because \(\alpha+2\beta\), \(\alpha-\beta\), and
\(2\alpha+\beta\) are not roots, this reduces to
\[
  p_\alpha(x)p_{\alpha+\beta}(Ax)
    =p_\alpha(Ax)p_{\alpha+\beta}(x).
\]
\SGARefLink{sga3:xxii:rem:4.1.3}{Exposé XXII, 4.1.3} now gives \(A=1\).

\SGARefTarget{sga3:XXIII:3.2.5}
\SGAParagraph{3.2.5}
Conjugating \eqref{sga3:XXIII:3.2:eqplus} by \(w_\alpha\), and using
\(A=1\), gives
\begin{equation}
  p_{\alpha+\beta}(zy)p_{-\alpha}(-x)
   =p_{-\alpha}(-x)p_{\alpha+\beta}(zy)
      p_\beta(-z^{-1}xy).
  \tag{+++}
  \label{sga3:XXIII:3.2:eqplusplusplus}
\end{equation}

\SGARefTarget{sga3:XXIII:3.2.6}
\SGAParagraph{3.2.6}
As before,
\[
  w_\alpha p_\beta(y)w_\alpha^{-1}=p_{\alpha+\beta}(zy).
\]
Since \(p_\alpha\) and \(p_{\alpha+\beta}\) commute,
\[
  p_\alpha(1)p_\beta(y)p_\alpha(-1)
    =p_{-\alpha}(1)p_{\alpha+\beta}(zy)p_{-\alpha}(-1).
\]
Using \eqref{sga3:XXIII:3.2:eqplus} and
\eqref{sga3:XXIII:3.2:eqplusplusplus}, this becomes
\[
  p_\beta(y)p_{\alpha+\beta}(-y)
    =p_{\alpha+\beta}(zy)p_\beta(-z^{-1}y).
\]
The two factors commute, and hence \(z=-1\).

\SGARefTarget{sga3:XXIII:3.2.7}
\SGAParagraph{3.2.7}
It remains to prove \SGARefLink{sga3:XXIII:3.2.1:item:i}{(i)}. One has
\[
  w_\alpha w_\beta w_\alpha
    =w_\alpha w_\beta w_\alpha^{-1}w_\alpha^2
    =w_{\alpha+\beta}^{-1}t_\alpha,
\]
and therefore
\[
\begin{aligned}
  w_\beta w_\alpha w_\beta w_\alpha w_\beta
    &=w_\beta w_{\alpha+\beta}^{-1}t_\alpha w_\beta\\
    &=w_\beta w_{\alpha+\beta}^{-1}w_\beta^{-1}
      \,s_\beta(t_\alpha)t_\beta\\
    &=w_\alpha t_\alpha t_\beta t_\beta
      =w_\alpha t_\alpha=w_\alpha^{-1}.
\end{aligned}
\]
Thus
\[
  (w_\alpha w_\beta)^3=(w_\beta w_\alpha)^3=e.
\]

\SGARefTarget{sga3:XXIII:3.2.8}
\SGAStatement{Remark 3.2.8}{Put
\(v_\alpha=f_\alpha(u_\alpha)\). Condition (v) of \SGARefLink{sga3:XXIII:2.3}{2.3} becomes
\[
\text{(A)}\quad
\begin{cases}
  \operatorname{int}(h_\alpha h_\beta h_\alpha)v_\beta=v_\beta,\\
  \operatorname{int}(h_\beta h_\alpha h_\beta)v_\alpha=v_\alpha,
\end{cases}
\]
and
\[
\text{(B)}\quad
\begin{cases}
  v_\beta v_\alpha
    =v_\alpha v_\beta\,(h_\beta v_\alpha h_\beta^{-1}),\\
  v_\beta(h_\beta v_\alpha h_\beta^{-1})
    =(h_\beta v_\alpha h_\beta^{-1})v_\beta,\\
  v_\alpha(h_\beta v_\alpha h_\beta^{-1})
    =(h_\beta v_\alpha h_\beta^{-1})v_\alpha.
\end{cases}
\]
Writing \(v_{\alpha+\beta}=\operatorname{int}(h_\beta)v_\alpha\),
the latter conditions are equivalently
\[
\text{(B)}\quad
\begin{cases}
  v_\beta v_\alpha=v_\alpha v_\beta v_{\alpha+\beta},\\
  v_\alpha v_{\alpha+\beta}=v_{\alpha+\beta}v_\alpha,\\
  v_\beta v_{\alpha+\beta}=v_{\alpha+\beta}v_\beta.
\end{cases}
\]}

\SGANextSectionTarget{sga3:xxiii:subsection:3-3}\SGASubsection{3.3}{Groups of Type \(B_2\)}

\SGARefTarget{sga3:XXIII:3.3.1}
\SGAInlineRefTarget{sga3:xxiii:proposition:3-3-1}\SGAStatement{Proposition 3.3.1}{Let \(S\) be a scheme and \(G\) a
pinned \(S\)-group of type \(B_2\). Write
\[
  \Delta=\{\alpha,\beta\},
  \qquad
  R^+=\{\alpha,\beta,\alpha+\beta,2\alpha+\beta\}.
\]
\begin{enumerate}
  \item[\SGAInlineRefTarget{sga3:xxiii:proposition:3-3-1:item:i}(i)]
    \[
      t_{\alpha\beta}=(w_\alpha w_\beta)^4
        =t_\alpha=(w_\beta w_\alpha)^4=t_{\beta\alpha}.
    \]
  \item[\SGAInlineRefTarget{sga3:XXIII:3.3.1:item:ii}\SGAInlineRefTarget{sga3:xxiii:proposition:3-3-1:item:ii}(ii)] If
    \[
      \operatorname{Ad}(w_\beta)X_\alpha=X_{\alpha+\beta},
      \qquad
      \operatorname{Ad}(w_\alpha)X_\beta=X_{2\alpha+\beta},
    \]
    then
    \[
    \begin{aligned}
      \operatorname{Ad}(w_\alpha)X_{\alpha+\beta}
        &=-X_{\alpha+\beta},&
      \operatorname{Ad}(w_\alpha)X_{2\alpha+\beta}&=X_\beta,\\
      \operatorname{Ad}(w_\beta)X_{\alpha+\beta}
        &=-X_\alpha,&
      \operatorname{Ad}(w_\beta)X_{2\alpha+\beta}&=X_{2\alpha+\beta}.
    \end{aligned}
    \]
  \item[\SGAInlineRefTarget{sga3:XXIII:3.3.1:item:iii}(iii)] Put
    \[
    \begin{aligned}
      p_{\alpha+\beta}(x)
        &=\exp(xX_{\alpha+\beta})
          =\operatorname{int}(w_\beta)p_\alpha(x),\\
      p_{2\alpha+\beta}(x)
        &=\exp(xX_{2\alpha+\beta})
          =\operatorname{int}(w_\alpha)p_\beta(x).
    \end{aligned}
    \]
    Then
    \[
    \begin{aligned}
      p_\beta(y)p_\alpha(x)
        &=p_\alpha(x)p_\beta(y)p_{\alpha+\beta}(xy)
          p_{2\alpha+\beta}(x^2y),\\
      p_{\alpha+\beta}(y)p_\alpha(x)
        &=p_\alpha(x)p_{\alpha+\beta}(y)
          p_{2\alpha+\beta}(2xy).
    \end{aligned}
    \]
\end{enumerate}}

\SGARefTarget{sga3:XXIII:3.3.2}
\SGAParagraph{3.3.2}
The proof occupies \SGARefLink{sga3:XXIII:3.3.2}{3.3.2}--\SGARefLink{sga3:XXIII:3.3.6}{3.3.6}. We have
\[
  (\beta^*,\alpha)=-1,\qquad(\alpha^*,\beta)=-2,
\]
so \(\alpha(t_\beta)=-1\) and \(\beta(t_\alpha)=1\). In particular,
\(\beta(t_\alpha)=\alpha(t_\alpha)=1\), so \(t_\alpha\) is a section
of \(\operatorname{Centr}(G)\). Put
\[
  \operatorname{Ad}(w_\beta)X_\alpha=X_{\alpha+\beta},
  \qquad
  \operatorname{Ad}(w_\alpha)X_\beta=X_{2\alpha+\beta}.
\]
Then
\[
\begin{aligned}
  \operatorname{Ad}(w_\beta)X_{\alpha+\beta}
    &=\alpha(t_\beta)X_\alpha=-X_\alpha,\\
  \operatorname{Ad}(w_\alpha)X_{2\alpha+\beta}
    &=\beta(t_\alpha)X_\beta=X_\beta,\\
  \operatorname{Ad}(w_\beta)X_{2\alpha+\beta}
    &=X_{2\alpha+\beta}.
\end{aligned}
\]
For some \(k\in\mathbf G_m(S)\),
\[
  \operatorname{Ad}(w_\alpha)X_{\alpha+\beta}
    =kX_{\alpha+\beta}.
\]
By \SGARefLink{sga3:xxii:cor:5.5.2}{Exposé XXII, 5.5.2}, there are
\(A,B,C\in\mathbf G_a(S)\) such that
\begin{align}
  p_\beta(y)p_\alpha(x)
    &=p_\alpha(x)p_\beta(y)p_{\alpha+\beta}(Axy)
      p_{2\alpha+\beta}(Bx^2y), \tag{1}\label{sga3:XXIII:3.3:eq1}\\
  p_{\alpha+\beta}(y)p_\alpha(x)
    &=p_\alpha(x)p_{\alpha+\beta}(y)
      p_{2\alpha+\beta}(Cxy). \tag{2}\label{sga3:XXIII:3.3:eq2}
\end{align}
It remains to prove
\[
  A=B=1,\qquad C=2,\qquad k=-1.
\]

\SGARefTarget{sga3:XXIII:3.3.3}
\SGAParagraph{3.3.3}
Conjugating \eqref{sga3:XXIII:3.3:eq1} by \(w_\alpha\) gives
\begin{equation}
\begin{aligned}
  p_{2\alpha+\beta}(y)p_{-\alpha}(-x)
    ={}&p_{-\alpha}(-x)p_{2\alpha+\beta}(y)\\
      &{}\cdot p_{\alpha+\beta}(Akxy)p_\beta(Bx^2y).
\end{aligned}
\tag{3}\label{sga3:XXIII:3.3:eq3}
\end{equation}
Transforming \eqref{sga3:XXIII:3.3:eq2} in the same way gives
\begin{equation}
  p_{\alpha+\beta}(ky)p_{-\alpha}(-x)
   =p_{-\alpha}(-x)p_{\alpha+\beta}(ky)p_\beta(Cxy).
\tag{4}\label{sga3:XXIII:3.3:eq4}
\end{equation}
Conjugating \eqref{sga3:XXIII:3.3:eq1} by \(w_\beta\) gives
\begin{equation}
\begin{aligned}
  p_{-\beta}(-y)p_{\alpha+\beta}(x)
    ={}&p_{\alpha+\beta}(x)p_{-\beta}(-y)\\
      &{}\cdot p_\alpha(-Axy)p_{2\alpha+\beta}(Bx^2y).
\end{aligned}
\tag{5}\label{sga3:XXIII:3.3:eq5}
\end{equation}

\SGARefTarget{sga3:XXIII:3.3.4}
\SGAParagraph{3.3.4}
Write
\[
  w_\beta p_\alpha(x)w_\beta^{-1}=p_{\alpha+\beta}(x).
\]
Because \(\alpha+2\beta\) is not a root,
\[
  p_\beta(1)p_\alpha(x)p_\beta(-1)
    =p_{-\beta}(1)p_{\alpha+\beta}(x)p_{-\beta}(-1).
\]
Using \eqref{sga3:XXIII:3.3:eq1} on the left and
\eqref{sga3:XXIII:3.3:eq5} on the right yields
\[
\begin{aligned}
  &p_\alpha(x)p_\beta(1)p_{\alpha+\beta}(Ax)
    p_{2\alpha+\beta}(Bx^2)p_\beta(-1)\\
  &\quad=p_{\alpha+\beta}(x)p_{-\beta}(1)p_\alpha(Ax)
    p_{2\alpha+\beta}(-Bx^2)p_{-\beta}(-1).
\end{aligned}
\]
After commuting the indicated factors and applying
\eqref{sga3:XXIII:3.3:eq2}, this becomes
\[
\begin{aligned}
  p_\alpha(x)p_{\alpha+\beta}(Ax)p_{2\alpha+\beta}(Bx^2)
  ={}&p_\alpha(Ax)p_{\alpha+\beta}(x)\\
     &{}\cdot p_{2\alpha+\beta}((AC-B)x^2).
\end{aligned}
\]
Hence
\[
  A=1,\qquad C=2B.
\]

\SGARefTarget{sga3:XXIII:3.3.5}
\SGAParagraph{3.3.5}
Now write
\[
  w_\alpha p_\beta(y)w_\alpha^{-1}=p_{2\alpha+\beta}(y).
\]
Since \(3\alpha+\beta\) is not a root,
\[
  p_\alpha(1)p_\beta(y)p_\alpha(-1)
    =p_{-\alpha}(1)p_{2\alpha+\beta}(y)p_{-\alpha}(-1).
\]
Using \eqref{sga3:XXIII:3.3:eq1} on the left and
\eqref{sga3:XXIII:3.3:eq3} on the right gives
\[
\begin{aligned}
  p_\beta(y)p_{\alpha+\beta}(-Ay)p_{2\alpha+\beta}(By)
  ={}&p_{2\alpha+\beta}(y)p_{\alpha+\beta}(Aky)p_\beta(By).
\end{aligned}
\]
The three groups commute, so
\[
  B=1,\qquad -A=Ak.
\]
Consequently
\[
  A=B=1,\qquad C=2,\qquad k=-1.
\]

\SGARefTarget{sga3:XXIII:3.3.6}
\SGAParagraph{3.3.6}
We have proved \SGARefLink{sga3:XXIII:3.3.1:item:ii}{(ii)} and \SGARefLink{sga3:XXIII:3.3.1:item:iii}{(iii)}. Since
\[
  s_\beta(t_\alpha)
    =t_\alpha t_\beta^{(\alpha^*,\beta)}
    =t_\alpha,
\]
one obtains successively
\[
\begin{aligned}
  w_\alpha w_\beta w_\alpha
    &=w_{2\alpha+\beta}t_\alpha,\\
  w_\beta w_\alpha w_\beta w_\alpha w_\beta
    &=w_{2\alpha+\beta}t_\alpha t_\beta,\\
  w_\alpha w_\beta w_\alpha w_\beta w_\alpha w_\beta w_\alpha
    &=w_\beta^{-1}t_\alpha.
\end{aligned}
\]
\NDE{The French prints
\((\alpha^*,\beta)=2\); the convention used immediately above gives
\(-2\). The conclusion is unchanged because \(t_\beta^2=e\).}
It follows that
\[
  (w_\beta w_\alpha)^4=t_\alpha,
  \qquad
  (w_\alpha w_\beta)^4=s_\beta(t_\alpha)=t_\alpha.
\]

\SGARefTarget{sga3:XXIII:3.3.7}
\SGAStatement{Remark 3.3.7}{Put
\[
  v_{\alpha+\beta}=\operatorname{int}(h_\beta)v_\alpha,
  \qquad
  v_{2\alpha+\beta}=\operatorname{int}(h_\alpha)v_\beta.
\]
Condition \SGARefLink{sga3:XXIII:2.3:item:v}{(v)} of \SGARefLink{sga3:XXIII:2.3}{2.3} becomes
\[
\text{(A)}\quad
\begin{cases}
  \operatorname{int}(h_\alpha h_\beta h_\alpha)v_\beta=v_\beta,\\
  \operatorname{int}(h_\beta h_\alpha h_\beta)v_\alpha=v_\alpha,
\end{cases}
\]
and
\[
\text{(B)}\quad
\begin{cases}
  v_\beta v_\alpha
    =v_\alpha v_\beta v_{\alpha+\beta}v_{2\alpha+\beta},\\
  v_{\alpha+\beta}v_\alpha
    =v_\alpha v_{\alpha+\beta}v_{2\alpha+\beta}^{\,2},\\
  v_{\alpha+\beta}v_\beta=v_\beta v_{\alpha+\beta},\\
  v_{2\alpha+\beta}v_\alpha=v_\alpha v_{2\alpha+\beta},\\
  v_{2\alpha+\beta}v_\beta=v_\beta v_{2\alpha+\beta},\\
  v_{2\alpha+\beta}v_{\alpha+\beta}
    =v_{\alpha+\beta}v_{2\alpha+\beta}.
\end{cases}
\]}

% END inputs/XXIII/tex/components/06_expose_XXIII_type_A2_and_B2.tex

% BEGIN inputs/XXIII/tex/components/07_expose_XXIII_type_G2.tex
\SGANextSectionTarget{sga3:xxiii:subsection:3-4}\SGASubsection{3.4}{Groups of Type \(G_2\)}

\SGARefTarget{sga3:XXIII:3.4.1}
\SGAInlineRefTarget{sga3:xxiii:proposition:3-4-1}\SGAStatement{Proposition 3.4.1}{Let \(S\) be a scheme and \(G\) a
pinned \(S\)-group of type \(G_2\). Write
\[
  \Delta=\{\alpha,\beta\},\qquad
  R^+=\{\alpha,\beta,\alpha+\beta,2\alpha+\beta,
                 3\alpha+\beta,3\alpha+2\beta\}.
\]
\begin{enumerate}
  \item[\SGAInlineRefTarget{sga3:xxiii:proposition:3-4-1:item:i}\SGAInlineRefTarget{sga3:XXIII:3.4.1:item:i}(i)]
    \[
      t_{\alpha\beta}=(w_\alpha w_\beta)^6
        =e=(w_\beta w_\alpha)^6=t_{\beta\alpha}.
    \]
  \item[\SGAInlineRefTarget{sga3:XXIII:3.4.1:item:ii}\SGAInlineRefTarget{sga3:xxiii:proposition:3-4-1:item:ii}(ii)] If
    \[
    \begin{aligned}
      \operatorname{Ad}(w_\beta)X_\alpha&=X_{\alpha+\beta},&
      \operatorname{Ad}(w_\alpha)X_{\alpha+\beta}&=X_{2\alpha+\beta},\\
      \operatorname{Ad}(w_\alpha)X_\beta&=-X_{3\alpha+\beta},&
      \operatorname{Ad}(w_\beta)X_{3\alpha+\beta}&=X_{3\alpha+2\beta},
    \end{aligned}
    \]
    then
    \[
    \begin{aligned}
      \operatorname{Ad}(w_\alpha)X_{2\alpha+\beta}
        &=-X_{\alpha+\beta},&
      \operatorname{Ad}(w_\alpha)X_{3\alpha+\beta}&=X_\beta,\\
      \operatorname{Ad}(w_\alpha)X_{3\alpha+2\beta}
        &=X_{3\alpha+2\beta},&
      \operatorname{Ad}(w_\beta)X_{\alpha+\beta}&=-X_\alpha,\\
      \operatorname{Ad}(w_\beta)X_{2\alpha+\beta}
        &=X_{2\alpha+\beta},&
      \operatorname{Ad}(w_\beta)X_{3\alpha+2\beta}&=-X_{3\alpha+\beta}.
    \end{aligned}
    \]
  \item[\SGAInlineRefTarget{sga3:XXIII:3.4.1:item:iii}(iii)] Put
    \[
    \begin{aligned}
      p_{\alpha+\beta}(x)
        &=\exp(xX_{\alpha+\beta})
          =\operatorname{int}(w_\beta)p_\alpha(x),\\
      p_{2\alpha+\beta}(x)
        &=\exp(xX_{2\alpha+\beta})
          =\operatorname{int}(w_\alpha w_\beta)p_\alpha(x),\\
      p_{3\alpha+\beta}(x)
        &=\exp(xX_{3\alpha+\beta})
          =\operatorname{int}(w_\alpha)p_\beta(-x),\\
      p_{3\alpha+2\beta}(x)
        &=\exp(xX_{3\alpha+2\beta})
          =\operatorname{int}(w_\beta w_\alpha)p_\beta(-x).
    \end{aligned}
    \]
    Then
    \[
    \begin{aligned}
      p_\beta(y)p_\alpha(x)
        ={}&p_\alpha(x)p_\beta(y)p_{\alpha+\beta}(xy)
          p_{2\alpha+\beta}(x^2y)p_{3\alpha+\beta}(x^3y)
          p_{3\alpha+2\beta}(x^3y^2),\\
      p_{\alpha+\beta}(y)p_\alpha(x)
        ={}&p_\alpha(x)p_{\alpha+\beta}(y)
          p_{2\alpha+\beta}(2xy)p_{3\alpha+\beta}(3x^2y)
          p_{3\alpha+2\beta}(3xy^2),\\
      p_{2\alpha+\beta}(y)p_\alpha(x)
        ={}&p_\alpha(x)p_{2\alpha+\beta}(y)
          p_{3\alpha+\beta}(3xy),\\
      p_{3\alpha+\beta}(y)p_\beta(x)
        ={}&p_\beta(x)p_{3\alpha+\beta}(y)
          p_{3\alpha+2\beta}(-xy),\\
      p_{2\alpha+\beta}(y)p_{\alpha+\beta}(x)
        ={}&p_{\alpha+\beta}(x)p_{2\alpha+\beta}(y)
          p_{3\alpha+2\beta}(3xy).
    \end{aligned}
    \]
\end{enumerate}}

\SGARefTarget{sga3:XXIII:3.4.2}
\SGAParagraph{3.4.2}
The proof occupies \SGARefLink{sga3:XXIII:3.4.2}{3.4.2}--\SGARefLink{sga3:XXIII:3.4.9}{3.4.9}. Since
\[
  (\beta^*,\alpha)=-1,\qquad(\alpha^*,\beta)=-3,
\]
we have \(\beta(t_\alpha)=\alpha(t_\beta)=-1\). Define the four
nonsimple root vectors as in \SGARefLink{sga3:XXIII:3.4.1:item:ii}{(ii)}. Then
\[
\begin{aligned}
  \operatorname{Ad}(w_\beta)X_{\alpha+\beta}
    &=\alpha(t_\beta)X_\alpha=-X_\alpha,\\
  \operatorname{Ad}(w_\alpha)X_{2\alpha+\beta}
    &=\alpha(t_\alpha)\beta(t_\alpha)X_{\alpha+\beta}
      =-X_{\alpha+\beta},\\
  \operatorname{Ad}(w_\alpha)X_{3\alpha+\beta}
    &=-\beta(t_\alpha)X_\beta=X_\beta,\\
  \operatorname{Ad}(w_\beta)X_{3\alpha+2\beta}
    &=\alpha(t_\beta)^3\beta(t_\beta)X_{3\alpha+\beta}
      =-X_{3\alpha+\beta}.
\end{aligned}
\]
Since \((3\alpha+2\beta)\pm\alpha\) and
\((2\alpha+\beta)\pm\beta\) are not roots,
\[
  \operatorname{Ad}(w_\alpha)X_{3\alpha+2\beta}=X_{3\alpha+2\beta},
  \qquad
  \operatorname{Ad}(w_\beta)X_{2\alpha+\beta}=X_{2\alpha+\beta}.
\]

\SGARefTarget{sga3:XXIII:3.4.3}
\SGAParagraph{3.4.3}
By \SGARefLink{sga3:xxii:cor:5.5.2}{Exposé XXII, 5.5.2}, there are sections
\[
  A,B,C,D,E,F,G,H,J\in\mathbf G_a(S)
\]
such that
\begin{align}
 p_\beta(y)p_\alpha(x)
  ={}&p_\alpha(x)p_\beta(y)p_{\alpha+\beta}(Axy)
     p_{2\alpha+\beta}(Bx^2y)p_{3\alpha+\beta}(Cx^3y)
     p_{3\alpha+2\beta}(Dx^3y^2), \tag{1}\label{sga3:XXIII:3.4:eq1}\\
 p_{\alpha+\beta}(y)p_\alpha(x)
  ={}&p_\alpha(x)p_{\alpha+\beta}(y)p_{2\alpha+\beta}(Exy)
     p_{3\alpha+\beta}(Fx^2y)p_{3\alpha+2\beta}(Gxy^2),
     \tag{2}\label{sga3:XXIII:3.4:eq2}\\
 p_{2\alpha+\beta}(y)p_\alpha(x)
  ={}&p_\alpha(x)p_{2\alpha+\beta}(y)p_{3\alpha+\beta}(Hxy),
     \tag{3}\label{sga3:XXIII:3.4:eq3}\\
 p_{3\alpha+\beta}(y)p_\beta(x)
  ={}&p_\beta(x)p_{3\alpha+\beta}(y)p_{3\alpha+2\beta}(Jxy).
     \tag{4}\label{sga3:XXIII:3.4:eq4}
\end{align}

\SGARefTarget{sga3:XXIII:3.4.4}
\SGAParagraph{3.4.4}
Conjugating (1), (3), and (4) by \(w_\beta\) gives
\begin{align}
 p_{-\beta}(-y)p_{\alpha+\beta}(x)
  ={}&p_{\alpha+\beta}(x)p_{-\beta}(-y)p_\alpha(-Axy)
     p_{2\alpha+\beta}(Bx^2y)\notag\\
    &{}\cdot p_{3\alpha+2\beta}(Cx^3y)
     p_{3\alpha+\beta}(-Dx^3y^2), \tag{5}\label{sga3:XXIII:3.4:eq5}\\
 p_{2\alpha+\beta}(y)p_{\alpha+\beta}(x)
  ={}&p_{\alpha+\beta}(x)p_{2\alpha+\beta}(y)
     p_{3\alpha+2\beta}(Hxy), \tag{6}\label{sga3:XXIII:3.4:eq6}\\
 p_{3\alpha+2\beta}(y)p_{-\beta}(-x)
  ={}&p_{-\beta}(-x)p_{3\alpha+2\beta}(y)
     p_{3\alpha+\beta}(-Jxy). \tag{7}\label{sga3:XXIII:3.4:eq7}
\end{align}
Conjugating (1) by \(w_\alpha\) gives
\begin{align}
 p_{3\alpha+\beta}(-y)p_{-\alpha}(-x)
  ={}&p_{-\alpha}(-x)p_{3\alpha+\beta}(-y)
     p_{2\alpha+\beta}(Axy)p_{\alpha+\beta}(-Bx^2y)\notag\\
    &{}\cdot p_\beta(Cx^3y)p_{3\alpha+\beta}(Dx^3y^2).
     \tag{8}\label{sga3:XXIII:3.4:eq8}
\end{align}

\SGARefTarget{sga3:XXIII:3.4.5}
\SGAParagraph{3.4.5}
The identity
\[
  w_\beta p_\alpha(x)w_\beta^{-1}=p_{\alpha+\beta}(x)
\]
may be written, since \(\alpha+2\beta\) is not a root,
\begin{equation}
  p_\beta(1)p_\alpha(x)p_\beta(-1)
   =p_{-\beta}(1)p_{\alpha+\beta}(x)p_{-\beta}(-1).
  \tag{9}\label{sga3:XXIII:3.4:eq9}
\end{equation}
Using (1) and then (4), its left side becomes
\begin{align}
 &p_\alpha(x)p_{\alpha+\beta}(Ax)p_{2\alpha+\beta}(Bx^2)
   p_{3\alpha+\beta}(Cx^3)
   p_{3\alpha+2\beta}((D-CJ)x^3). \tag{10}\label{sga3:XXIII:3.4:eq10}
\end{align}
Using (5) and then \SGARefLink{sga3:XXIII:3.4:eq7}{(7)}, its right side becomes
\[
 p_{\alpha+\beta}(x)p_\alpha(Ax)p_{2\alpha+\beta}(-Bx^2)
 p_{3\alpha+2\beta}(-Cx^3)p_{3\alpha+\beta}((CJ-D)x^3).
\]
Applying (2) transforms this into
\begin{align}
 &p_\alpha(Ax)p_{\alpha+\beta}(x)
 p_{2\alpha+\beta}((AE-B)x^2)
 p_{3\alpha+\beta}((A^2F+CJ-D)x^3)\notag\\
 &\hspace{35mm}\cdot
 p_{3\alpha+2\beta}((AG-C)x^3). \tag{12}\label{sga3:XXIII:3.4:eq12}
\end{align}
Comparing \SGARefLink{sga3:XXIII:3.4:eq10}{(10)} and \SGARefLink{sga3:XXIII:3.4:eq12}{(12)} gives
\[
  A=1,\qquad E=2B,\qquad C+D=F+CJ,\qquad F=G.
\]

\SGARefTarget{sga3:XXIII:3.4.6}
\SGAParagraph{3.4.6}
Write
\[
  w_\alpha p_\beta(y)w_\alpha^{-1}=p_{3\alpha+\beta}(-y).
\]
Since \(4\alpha+\beta\) is not a root,
\begin{equation}
 p_\alpha(1)p_\beta(y)p_\alpha(-1)
  =p_{-\alpha}(1)p_{3\alpha+\beta}(-y)p_{-\alpha}(-1).
 \tag{13}\label{sga3:XXIII:3.4:eq13}
\end{equation}
Using (1), the left side is
\[
 p_\beta(y)p_{\alpha+\beta}(-Ay)p_{2\alpha+\beta}(By)
 p_{3\alpha+\beta}(-Cy)p_{3\alpha+2\beta}(-Dy^2).
\]
Using \SGARefLink{sga3:XXIII:3.4:eq8}{(8)}, (6), and (4), the right side is
\[
\begin{aligned}
 p_\beta(Cy)p_{\alpha+\beta}(-By)p_{2\alpha+\beta}(Ay)
 p_{3\alpha+\beta}(-y)
 p_{3\alpha+2\beta}((D-CJ-ABH)y^2).
\end{aligned}
\]
Thus
\[
  C=1,\qquad A=B,\qquad D-CJ-ABH=-D.
\]
Together with the preceding identities, this yields
\[
 A=B=C=1,\quad E=2,\quad F=G,\quad
 D+1=F+J,\quad 2D=H+J.
\]

\SGARefTarget{sga3:XXIII:3.4.7}
\SGAParagraph{3.4.7}
From
\[
  w_\beta p_{3\alpha+\beta}(x)w_\beta^{-1}
    =p_{3\alpha+2\beta}(x)
\]
one obtains
\[
  p_{3\alpha+\beta}(x)p_{3\alpha+2\beta}(-Jx)
   =p_{3\alpha+2\beta}(x)p_{3\alpha+\beta}(-Jx),
\]
and therefore \(J=-1\).

\SGARefTarget{sga3:XXIII:3.4.8}
\SGAParagraph{3.4.8}
Finally,
\[
  w_\alpha p_{\alpha+\beta}(y)w_\alpha^{-1}=p_{2\alpha+\beta}(y),
\]
or equivalently
\[
 p_\alpha(1)p_{\alpha+\beta}(y)p_\alpha(-1)
  =p_{-\alpha}(1)p_\alpha(-1)p_{2\alpha+\beta}(y)
    p_\alpha(1)p_{-\alpha}(-1).
\]
Using (2) on the left and (3) on the right gives
\[
\begin{aligned}
 &p_{\alpha+\beta}(y)p_{2\alpha+\beta}(-Ey)
   p_{3\alpha+\beta}(Fy)p_{3\alpha+2\beta}(-Gy^2)\\
 &\qquad=p_{-\alpha}(1)p_{2\alpha+\beta}(y)
   p_{3\alpha+\beta}(Hy)p_{-\alpha}(-1).
\end{aligned}
\]
Commuting \(p_{-\alpha}(-1)\) first past
\(p_{3\alpha+\beta}(Hy)\), and then past \(p_{2\alpha+\beta}(y)\),
introduces no new \(p_{3\alpha+\beta}\)-factor. Comparing those
factors gives \(F=H\). Hence
\[
  2D=D+1,\qquad D=1,\qquad F=G=H=2D-J=3.
\]
This determines all the coefficients and proves \SGARefLink{sga3:XXIII:3.4.1:item:iii}{(iii)}.

\SGARefTarget{sga3:XXIII:3.4.9}
\SGAParagraph{3.4.9}
For \SGARefLink{sga3:XXIII:3.4.1:item:i}{(i)}, one obtains successively
\[
\begin{aligned}
 w_\alpha w_\beta w_\alpha
  &=w_{3\alpha+\beta}^{-1}t_\alpha,\\
 w_\beta(w_\alpha w_\beta)^2
  &=w_{3\alpha+2\beta}^{-1}t_\alpha,\\
 w_\alpha(w_\beta w_\alpha)^3
  &=w_{3\alpha+2\beta}^{-1},\\
 w_\beta(w_\alpha w_\beta)^4
  &=w_{3\alpha+\beta}t_\beta,\\
 w_\alpha(w_\beta w_\alpha)^5
  &=w_\beta^{-1}.
\end{aligned}
\]
Consequently
\[
  (w_\alpha w_\beta)^6=(w_\beta w_\alpha)^6=e.
\]

\SGARefTarget{sga3:XXIII:3.4.10}
\SGAStatement{Remark 3.4.10}{Condition \SGARefLink{sga3:XXIII:2.3:item:v}{(v)} of \SGARefLink{sga3:XXIII:2.3}{2.3} consists first of
\[
\text{(A)}\quad
\begin{cases}
 \operatorname{int}(h_\alpha h_\beta h_\alpha h_\beta h_\alpha)v_\beta
   =v_\beta,\\
 \operatorname{int}(h_\beta h_\alpha h_\beta h_\alpha h_\beta)v_\alpha
   =v_\alpha,
\end{cases}
\]
and, after setting
\[
\begin{aligned}
 v_{\alpha+\beta}&=\operatorname{int}(h_\beta)v_\alpha,&
 v_{2\alpha+\beta}&=\operatorname{int}(h_\alpha h_\beta)v_\alpha,\\
 v_{3\alpha+\beta}&=\operatorname{int}(h_\alpha)v_\beta^{-1},&
 v_{3\alpha+2\beta}&=\operatorname{int}(h_\beta h_\alpha)v_\beta^{-1},
\end{aligned}
\]
of the commutation relations
\[
\text{(B)}\quad
\begin{aligned}
 v_\beta v_\alpha
  &=v_\alpha v_\beta v_{\alpha+\beta}v_{2\alpha+\beta}
    v_{3\alpha+\beta}v_{3\alpha+2\beta},\\
 v_{\alpha+\beta}v_\alpha
  &=v_\alpha v_{\alpha+\beta}v_{2\alpha+\beta}^{\,2}
    v_{3\alpha+\beta}^{\,3}v_{3\alpha+2\beta}^{\,3},\\
 v_{2\alpha+\beta}v_\alpha
  &=v_\alpha v_{2\alpha+\beta}v_{3\alpha+\beta}^{\,3},\\
 v_{3\alpha+\beta}v_\alpha&=v_\alpha v_{3\alpha+\beta},&
 v_{3\alpha+2\beta}v_\alpha&=v_\alpha v_{3\alpha+2\beta},\\
 v_{\alpha+\beta}v_\beta&=v_\beta v_{\alpha+\beta},&
 v_{2\alpha+\beta}v_\beta&=v_\beta v_{2\alpha+\beta},\\
 v_{3\alpha+\beta}v_\beta
  &=v_\beta v_{3\alpha+\beta}v_{3\alpha+2\beta}^{-1},\\
 v_{3\alpha+2\beta}v_\beta&=v_\beta v_{3\alpha+2\beta},\\
 v_{2\alpha+\beta}v_{\alpha+\beta}
  &=v_{\alpha+\beta}v_{2\alpha+\beta}v_{3\alpha+2\beta}^{\,3},\\
 v_{3\alpha+\beta}v_{\alpha+\beta}
  &=v_{\alpha+\beta}v_{3\alpha+\beta},\\
 v_{3\alpha+2\beta}v_{\alpha+\beta}
  &=v_{\alpha+\beta}v_{3\alpha+2\beta},\\
 v_{3\alpha+\beta}v_{2\alpha+\beta}
  &=v_{2\alpha+\beta}v_{3\alpha+\beta},\\
 v_{3\alpha+2\beta}v_{2\alpha+\beta}
  &=v_{2\alpha+\beta}v_{3\alpha+2\beta},\\
 v_{3\alpha+2\beta}v_{3\alpha+\beta}
  &=v_{3\alpha+\beta}v_{3\alpha+2\beta}.
\end{aligned}
\]}

% END inputs/XXIII/tex/components/07_expose_XXIII_type_G2.tex

% BEGIN inputs/XXIII/tex/components/08_expose_XXIII_generators_relations_explicit.tex
\SGANextSectionTarget{sga3:xxiii:subsection:3-5}\SGASubsection{3.5}{Explicit Form of the Generators-and-Relations
Theorem}

\SGARefTarget{sga3:XXIII:3.5.1}
\SGAStatement{Theorem 3.5.1}{Let \(S\) be a scheme and \(G\) a
pinned \(S\)-group, with maximal torus \(T\), system of simple roots
\(\Delta\), and elements
\[
  u_\alpha\in U_\alpha(S)^\times,\qquad
  w_\alpha\in\operatorname{Norm}_G(T)(S)
  \qquad(\alpha\in\Delta)
\]
defined by the pinning. Let \(H\) be an fppf \(S\)-sheaf in groups,
let
\[
  f_T:T\to H,\qquad f_\alpha:U_\alpha\to H
  \quad(\alpha\in\Delta)
\]
be morphisms of groups, and let \(h_\alpha\in H(S)\). Put
\[
  v_\alpha=f_\alpha(u_\alpha).
\]
There exists a morphism of groups \(f:G\to H\) extending \(f_T\) and
all the \(f_\alpha\), and satisfying
\(f(w_\alpha)=h_\alpha\), if and only if:
\begin{enumerate}
  \item[(i)] for every \(S'\to S\), \(\alpha\in\Delta\),
    \(t\in T(S')\), and \(x\in U_\alpha(S')\),
    \begin{equation}
      \operatorname{int}(f_T(t))f_\alpha(x)
       =f_\alpha(\operatorname{int}(t)x)
       =f_\alpha(x^{\alpha(t)}); \tag{1}\label{sga3:xxiii:theorem:3-5-1:equation:1}
    \end{equation}
  \item[(ii)] for every \(\alpha\in\Delta\), every \(S'\to S\), and
    every \(t\in T(S')\),
    \begin{equation}
      \operatorname{int}(h_\alpha)f_T(t)
       =f_T(s_\alpha(t))
       =f_T\!\left(t\alpha^*(\alpha(t)^{-1})\right); \tag{2}\label{sga3:xxiii:theorem:3-5-1:equation:2}
    \end{equation}
  \item[(iii)] for every \(\alpha\in\Delta\),
    \begin{equation}
      h_\alpha^2=f_T(\alpha^*(-1)), \tag{\(3_1\)}\label{sga3:xxiii:theorem:3-5-1:equation:3-1}
    \end{equation}
    and
    \begin{equation}
      (h_\alpha v_\alpha)^3=e; \tag{4}\label{sga3:xxiii:theorem:3-5-1:equation:4}
    \end{equation}
  \item[\SGAInlineRefTarget{sga3:XXIII:3.5.1:item:iv}(iv)] for distinct \(\alpha,\beta\in\Delta\):
    \begin{enumerate}
      \item[(a)] the appropriate braid relation is
        \begin{equation}
        \tag{\(3_2\)}\label{sga3:xxiii:theorem:3-5-1:equation:3-2}
        \begin{cases}
        (h_\alpha h_\beta)^2
          =f_T(\alpha^*(-1)\beta^*(-1)),
          &(\alpha^*,\beta)=0,\\
        (h_\alpha h_\beta)^3=e,
          &(\alpha^*,\beta)=-1,\\
        (h_\alpha h_\beta)^4=f_T(\alpha^*(-1)),
          &(\alpha^*,\beta)=-2,\\
        (h_\alpha h_\beta)^6=e,
          &(\alpha^*,\beta)=-3;
        \end{cases}
        \end{equation}
      \item[(b)] the corresponding relations (A) and (B) of
        \SGARefLink{sga3:XXIII:3.1.3}{3.1.3}, \SGARefLink{sga3:XXIII:3.2.8}{3.2.8}, \SGARefLink{sga3:XXIII:3.3.7}{3.3.7}, or \SGARefLink{sga3:XXIII:3.4.10}{3.4.10} hold.
    \end{enumerate}
\end{enumerate}
When it exists, \(f\) is necessarily unique.}

This follows immediately from \SGARefLink{sga3:XXIII:2.3}{Theorem 2.3}, \SGARefLink{sga3:XXIII:2.6}{Remark 2.6}, and the
calculations in each rank-two case.

\SGARefTarget{sga3:XXIII:3.5.2}
\SGAStatement{Remark 3.5.2}{The result may equivalently be presented
by giving morphisms
\[
  a_\alpha:T\!\cdot U_\alpha\longrightarrow H,
  \qquad
  b_\alpha:\operatorname{Norm}_{Z_\alpha}(T)\longrightarrow H
  \qquad(\alpha\in\Delta),
\]
and putting
\[
  h_\alpha=b_\alpha(w_\alpha),
  \qquad
  v_\alpha=a_\alpha(u_\alpha).
\]
The required conditions are:
\begin{enumerate}
  \item[(1)] all the \(a_\alpha\) and \(b_\alpha\) have the same
    restriction to \(T\);
  \item[(2)] condition (4) and all the conditions in \SGARefLink{sga3:XXIII:3.5.1:item:iv}{(iv)} of
    \SGARefLink{sga3:XXIII:3.5.1}{Theorem 3.5.1} hold.
\end{enumerate}}

\SGARefTarget{sga3:XXIII:3.5.3}
\SGAParagraph{3.5.3}
Further applications will be given in the next Exposé. Here is one.
\SGARefLink{sga3:XXIII:3.5.1}{Theorem 3.5.1} describes \(G\) by generators and relations in the
category of fppf \(S\)-sheaves. More explicitly, for every
\(S'\to S\), let \(H(S')\) be the group generated by
\[
  T(S'),\qquad U_\alpha(S')\ (\alpha\in R),\qquad
  w_\alpha\ (\alpha\in R),
\]
subject to relations analogous to
\[
  (1),\ (2),\ (3_1),\ (4),\ (3_2),\ (A),\ (B).
\]
Then \(G\) is the sheaf associated with the presheaf
\(S'\mapsto H(S')\).

In particular, if \(S'\) is the spectrum of an algebraically closed
field \(k\), then \(G(S')=H(S')\). This follows immediately from the
Nullstellensatz in the form: an fppf-covering sieve over an
algebraically closed field is trivial. Thus \SGARefLink{sga3:XXIII:3.5.1}{Theorem 3.5.1} gives an
explicit presentation by generators and relations of the abstract
group \(G(k)\).

% END inputs/XXIII/tex/components/08_expose_XXIII_generators_relations_explicit.tex

% BEGIN inputs/XXIII/tex/components/09_expose_XXIII_fundamental_theorem.tex
\SGANextSectionTarget{sga3:xxiii:section:4}\SGASection{4}{Uniqueness of Pinned Groups: Fundamental Theorem}

\SGARefTarget{sga3:XXIII:4.1}
\SGAInlineRefTarget{sga3:xxiii:theorem:4-1}\SGAStatement{Theorem 4.1}{Let \(S\) be a nonempty scheme. The
functor \(R\) of \SGARefLink{sga3:XXIII:1.6}{1.6} is fully faithful. More precisely, let \(G,G'\)
be pinned \(S\)-groups, let \(p>0\) be an integer such that
\(x\mapsto x^p\) is an endomorphism of \(\mathbf G_{a,S}\), and let
\[
  h:R(G')\longrightarrow R(G)
\]
be a \(p\)-morphism of pinned root data. There is a unique morphism
of pinned groups \(f:G\to G'\) such that \(R(f)=h\).}

Uniqueness was proved in \SGARefLink{sga3:XXIII:1.9}{1.9}, so only existence remains. By
hypothesis there are a bijection \(d:R\xrightarrow{\sim}R'\) and a
map
\[
  q:R\longrightarrow\{p^n\mid n\in\mathbf N\}
\]
such that, for every \(\alpha\in R\),
\[
  h(d(\alpha))=q(\alpha)\alpha,
  \qquad
  {}^th(\alpha^*)=q(\alpha)d(\alpha)^*.
\]
In particular, \(G\) and \(G'\) have the same semisimple rank.

\SGARefTarget{sga3:XXIII:4.1.1}
\SGAParagraph{4.1.1}
Suppose that both semisimple ranks are \(0\). Then \(G,G'\) are
tori:
\[
  G=T=D_S(M),\qquad G'=T'=D_S(M'),
\]
and \(h:M'\to M\) is an ordinary homomorphism. Take \(f=D_S(h)\).

\SGARefTarget{sga3:XXIII:4.1.2}
\SGAParagraph{4.1.2}
Suppose that both semisimple ranks are \(1\), and put
\[
  f_T=D_S(h):T\longrightarrow T'.
\]
Writing \(\alpha'=d(\alpha)\), the diagram
\[
\begin{tikzcd}[column sep=large,row sep=large]
\mathbf G_{m,S}\ar[r,"\alpha^*"]\ar[d,"q(\alpha)"]
  &T\ar[r,"\alpha"]\ar[d,"f_T"']
  &\mathbf G_{m,S}\ar[d,"q(\alpha)"]\\
\mathbf G_{m,S}\ar[r,"\alpha'^*"]
  &T'\ar[r,"\alpha'"]
  &\mathbf G_{m,S}
\end{tikzcd}
\]
commutes. Apply \SGARefLink{sga3:xx:4.1}{Exposé XX, 4.1}.

\SGARefTarget{sga3:XXIII:4.1.3}
\SGAParagraph{4.1.3}
Suppose that both semisimple ranks are \(2\). \SGARefLink{sga3:XXI:7.5.3}{Exposé XXI, 7.5.3}
lists every possible \(h\); we verify the conditions of \SGARefLink{sga3:XXIII:2.5}{Corollary 2.5}
in each case. Write
\[
  \Delta=\{\alpha,\beta\},\qquad
  \Delta'=\{\alpha',\beta'\},\qquad
  d(\alpha)=\alpha',\quad d(\beta)=\beta'.
\]

\SGARefTarget{sga3:XXIII:4.1.4}
\SGAParagraph{4.1.4. \(G\) and \(G'\) of type
\(A_1\times A_1\)}
Here
\[
  h(\alpha')=q\alpha,\qquad h(\beta')=q_1\beta.
\]
Conditions \SGARefLink{sga3:XXIII:2.5:item:ii}{2.5(ii)} and (iii) follow from \SGARefLink{sga3:xxiii:proposition:3-1-2:item:ii}{Proposition 3.1.2(ii)},(iii).
For (i),
\[
\begin{aligned}
 D_S(h)t_{\alpha\beta}
  &=D_S(h)(t_\alpha t_\beta)\\
  &={}^th(\alpha^*)(-1)\,{}^th(\beta^*)(-1)\\
  &=\alpha'^*((-1)^q)\,\beta'^*((-1)^{q_1}).
\end{aligned}
\]
The assumption on \(x\mapsto x^p\), with \(S\ne\varnothing\), implies
that \(p=1\) or \(S\) has characteristic \(p\). Thus
\((-1)^q=-1\): if \(q\) is even, then \(p=2\) and \(1=-1\).
Consequently
\[
  D_S(h)t_{\alpha\beta}
   =\alpha'^*(-1)\beta'^*(-1)=t'_{\alpha'\beta'}.
\]

\SGARefTarget{sga3:XXIII:4.1.5}
\SGAParagraph{4.1.5. \(G\) and \(G'\) of type \(A_2\)}
Here
\[
  h(\alpha')=q\alpha,\qquad h(\beta')=q\beta.
\]
Put
\[
  X_{\alpha+\beta}=\operatorname{Ad}(w_\beta)X_\alpha,
  \qquad
  X'_{\alpha'+\beta'}=\operatorname{Ad}(w'_{\beta'})X'_{\alpha'}.
\]
Condition (i) follows as above from \SGARefLink{sga3:xxiii:proposition:3-2-1:item:i}{Proposition 3.2.1(i)}, and (ii)
follows immediately from \SGARefLink{sga3:xxiii:proposition:3-2-1:item:ii}{Proposition 3.2.1(ii)}. For (iii), the map
\[
\begin{aligned}
 p_\alpha(x)p_\beta(y)p_{\alpha+\beta}(z)
 \longmapsto
 p'_{\alpha'}(x^q)p'_{\beta'}(y^q)p'_{\alpha'+\beta'}(z^q)
\end{aligned}
\]
must be a homomorphism. The only nontrivial relation is
\[
 p_\beta(y)p_\alpha(x)
  =p_\alpha(x)p_\beta(y)p_{\alpha+\beta}(xy),
\]
which maps to
\[
 p'_{\beta'}(y^q)p'_{\alpha'}(x^q)
  =p'_{\alpha'}(x^q)p'_{\beta'}(y^q)
    p'_{\alpha'+\beta'}(x^qy^q).
\]

\SGARefTarget{sga3:XXIII:4.1.6}
\SGAParagraph{4.1.6}
The same argument treats types \(B_2\) and \(G_2\) when the
radicial exponents are equal, using \SGARefLink{sga3:xxiii:proposition:3-3-1}{Propositions 3.3.1} and \SGARefLink{sga3:xxiii:proposition:3-4-1}{3.4.1}.
It remains to treat the two exceptional cases in \SGARefLink{sga3:XXI:7.5.3}{Exposé XXI, 7.5.3}.

\SGARefTarget{sga3:XXIII:4.1.7}
\SGAParagraph{4.1.7. Type \(B_2\), characteristic \(2\),
\(q(\alpha)=2q\), \(q(\beta)=q\)}
The positive roots are
\[
 \{\alpha,\beta,\alpha+\beta,2\alpha+\beta\}
\]
and
\[
 \{\alpha',\beta',\alpha'+\beta',\alpha'+2\beta'\},
\]
where \(\alpha\) and \(\beta'\) are the short simple roots. We have
\[
\begin{aligned}
 h(\alpha')&=2q\alpha,&h(\beta')&=q\beta,\\
 h(\alpha'+\beta')&=q(2\alpha+\beta),&
 h(\alpha'+2\beta')&=2q(\alpha+\beta),
\end{aligned}
\]
so
\[
\begin{aligned}
 d(\alpha+\beta)&=\alpha'+2\beta',&q(\alpha+\beta)&=2q,\\
 d(2\alpha+\beta)&=\alpha'+\beta',&q(2\alpha+\beta)&=q.
\end{aligned}
\]
Put
\[
\begin{aligned}
 X_{\alpha+\beta}&=\operatorname{Ad}(w_\beta)X_\alpha,&
 X_{2\alpha+\beta}&=\operatorname{Ad}(w_\alpha)X_\beta,\\
 X'_{\alpha'+\beta'}&=\operatorname{Ad}(w'_{\alpha'})X'_{\beta'},&
 X'_{\alpha'+2\beta'}&=\operatorname{Ad}(w'_{\beta'})X'_{\alpha'}.
\end{aligned}
\]

For (i), characteristic \(2\) gives \(-1=1\), hence
\[
 t_{\alpha\beta}=t_\alpha=\alpha^*(-1)=e
   =\beta'^*(-1)=t'_{\beta'}=t'_{\alpha'\beta'}.
\]
For (ii), the two defining identities are
\[
\begin{aligned}
 \operatorname{Ad}(w_\alpha)X_\beta
  &=X_{2\alpha+\beta},&
 \operatorname{Ad}(w'_{\alpha'})X'_{\beta'}
  &=X'_{\alpha'+\beta'}=X'_{d(2\alpha+\beta)},\\
 \operatorname{Ad}(w_\beta)X_\alpha
  &=X_{\alpha+\beta},&
 \operatorname{Ad}(w'_{\beta'})X'_{\alpha'}
  &=X'_{\alpha'+2\beta'}=X'_{d(\alpha+\beta)}.
\end{aligned}
\]
\SGARefLink{sga3:xxiii:proposition:3-3-1:item:ii}{Proposition 3.3.1(ii)}, together with \(-1=1\), gives the four
remaining identities:
\[
\begin{aligned}
 \operatorname{Ad}(w_\alpha)X_{\alpha+\beta}&=X_{\alpha+\beta},&
 \operatorname{Ad}(w'_{\alpha'})X'_{d(\alpha+\beta)}
   &=X'_{d(\alpha+\beta)},\\
 \operatorname{Ad}(w_\alpha)X_{2\alpha+\beta}&=X_\beta,&
 \operatorname{Ad}(w'_{\alpha'})X'_{d(2\alpha+\beta)}
   &=X'_{d(\beta)},\\
 \operatorname{Ad}(w_\beta)X_{\alpha+\beta}&=X_\alpha,&
 \operatorname{Ad}(w'_{\beta'})X'_{d(\alpha+\beta)}
   &=X'_{d(\alpha)},\\
 \operatorname{Ad}(w_\beta)X_{2\alpha+\beta}&=X_{2\alpha+\beta},&
 \operatorname{Ad}(w'_{\beta'})X'_{d(2\alpha+\beta)}
   &=X'_{d(2\alpha+\beta)}.
\end{aligned}
\]

For (iii), the sole nontrivial relations are
\[
\begin{aligned}
 p_\beta(y)p_\alpha(x)
  &=p_\alpha(x)p_\beta(y)p_{\alpha+\beta}(xy)
    p_{2\alpha+\beta}(x^2y),\\
 p'_{\alpha'}(y')p'_{\beta'}(x')
  &=p'_{\beta'}(x')p'_{\alpha'}(y')
    p'_{\alpha'+\beta'}(x'y')
    p'_{\alpha'+2\beta'}(x'y'^2).
\end{aligned}
\]
The proposed map is
\[
\begin{aligned}
 &p_\alpha(x)p_\beta(y)p_{\alpha+\beta}(z)p_{2\alpha+\beta}(t)\\
 &\quad\longmapsto
 p'_{\alpha'}(x^{2q})p'_{\beta'}(y^q)
 p'_{\alpha'+2\beta'}(z^{2q})p'_{\alpha'+\beta'}(t^q).
\end{aligned}
\]
It respects the first relation because the second, with
\(y'=x^{2q}\) and \(x'=y^q\), reads
\[
\begin{aligned}
 p'_{\beta'}(y^q)p'_{\alpha'}(x^{2q})
  ={}&p'_{\alpha'}(x^{2q})p'_{\beta'}(y^q)
   p'_{\alpha'+2\beta'}((xy)^{2q})\\
   &{}\cdot p'_{\alpha'+\beta'}((x^2y)^q).
\end{aligned}
\]

\SGARefTarget{sga3:XXIII:4.1.8}
\SGAParagraph{4.1.8. Type \(G_2\), characteristic \(3\),
\(q(\alpha)=3q\), \(q(\beta)=q\)}
The positive roots are
\[
\{\alpha,\beta,\alpha+\beta,2\alpha+\beta,3\alpha+\beta,
  3\alpha+2\beta\}
\]
and
\[
\{\alpha',\beta',\alpha'+\beta',\alpha'+2\beta',
  \alpha'+3\beta',2\alpha'+3\beta'\},
\]
where again \(\alpha\) and \(\beta'\) are short. We have
\[
\begin{aligned}
 h(\alpha')&=3q\alpha,&h(\beta')&=q\beta,\\
 h(\alpha'+\beta')&=q(3\alpha+\beta),&
 h(\alpha'+2\beta')&=q(3\alpha+2\beta),\\
 h(\alpha'+3\beta')&=3q(\alpha+\beta),&
 h(2\alpha'+3\beta')&=3q(2\alpha+\beta),
\end{aligned}
\]
and hence
\[
\begin{aligned}
 d(\alpha+\beta)&=\alpha'+3\beta',&q(\alpha+\beta)&=3q,\\
 d(2\alpha+\beta)&=2\alpha'+3\beta',&q(2\alpha+\beta)&=3q,\\
 d(3\alpha+\beta)&=\alpha'+\beta',&q(3\alpha+\beta)&=q,\\
 d(3\alpha+2\beta)&=\alpha'+2\beta',&q(3\alpha+2\beta)&=q.
\end{aligned}
\]
Define
\[
\begin{aligned}
 X_{\alpha+\beta}&=\operatorname{Ad}(w_\beta)X_\alpha,&
 X_{2\alpha+\beta}&=\operatorname{Ad}(w_\alpha)X_{\alpha+\beta},\\
 X_{3\alpha+\beta}&=-\operatorname{Ad}(w_\alpha)X_\beta,&
 X_{3\alpha+2\beta}&=\operatorname{Ad}(w_\beta)X_{3\alpha+\beta},\\
 X'_{\alpha'+\beta'}&=-\operatorname{Ad}(w'_{\alpha'})X'_{\beta'},&
 X'_{\alpha'+2\beta'}&=\operatorname{Ad}(w'_{\beta'})X'_{\alpha'+\beta'},\\
 X'_{\alpha'+3\beta'}&=\operatorname{Ad}(w'_{\beta'})X'_{\alpha'},&
 X'_{2\alpha'+3\beta'}&=\operatorname{Ad}(w'_{\alpha'})X'_{\alpha'+3\beta'}.
\end{aligned}
\]
Condition (i) follows from \SGARefLink{sga3:xxiii:proposition:3-4-1:item:i}{Proposition 3.4.1(i)}. Four instances of
(ii) hold by definition:
\[
\begin{aligned}
 \operatorname{Ad}(w_\beta)X_\alpha&=X_{\alpha+\beta},&
 \operatorname{Ad}(w'_{\beta'})X'_{\alpha'}&=X'_{d(\alpha+\beta)},\\
 \operatorname{Ad}(w_\alpha)X_\beta&=-X_{3\alpha+\beta},&
 \operatorname{Ad}(w'_{\alpha'})X'_{\beta'}&=-X'_{d(3\alpha+\beta)},\\
 \operatorname{Ad}(w_\beta)X_{3\alpha+\beta}&=X_{3\alpha+2\beta},&
 \operatorname{Ad}(w'_{\beta'})X'_{d(3\alpha+\beta)}
   &=X'_{d(3\alpha+2\beta)},\\
 \operatorname{Ad}(w_\alpha)X_{\alpha+\beta}&=X_{2\alpha+\beta},&
 \operatorname{Ad}(w'_{\alpha'})X'_{d(\alpha+\beta)}
   &=X'_{d(2\alpha+\beta)}.
\end{aligned}
\]
The remaining instances follow from \SGARefLink{sga3:xxiii:proposition:3-4-1:item:ii}{Proposition 3.4.1(ii)}; for
example,
\[
\begin{aligned}
 \operatorname{Ad}(w_\alpha)X_{2\alpha+\beta}&=-X_{\alpha+\beta},&
 \operatorname{Ad}(w'_{\alpha'})X'_{d(2\alpha+\beta)}
   &=-X'_{d(\alpha+\beta)},\\
 \operatorname{Ad}(w_\beta)X_{3\alpha+2\beta}&=-X_{3\alpha+\beta},&
 \operatorname{Ad}(w'_{\beta'})X'_{d(3\alpha+2\beta)}
   &=-X'_{d(3\alpha+\beta)}.
\end{aligned}
\]

For (iii), in characteristic \(3\), the nontrivial relations in
\(U\) are
\[
\begin{aligned}
 p_\beta(y)p_\alpha(x)
  ={}&p_\alpha(x)p_\beta(y)p_{\alpha+\beta}(xy)
  p_{2\alpha+\beta}(x^2y)p_{3\alpha+\beta}(x^3y)
  p_{3\alpha+2\beta}(x^3y^2),\\
 p_{\alpha+\beta}(y)p_\alpha(x)
  ={}&p_\alpha(x)p_{\alpha+\beta}(y)p_{2\alpha+\beta}(-xy),\\
 p_{3\alpha+\beta}(y)p_\beta(x)
  ={}&p_\beta(x)p_{3\alpha+\beta}(y)p_{3\alpha+2\beta}(-xy).
\end{aligned}
\]
The corresponding relations in \(U'\) are
\[
\begin{aligned}
 p'_{\alpha'}(y')p'_{\beta'}(x')
  ={}&p'_{\beta'}(x')p'_{\alpha'}(y')
  p'_{\alpha'+\beta'}(-x'y')p'_{\alpha'+2\beta'}(-x'^2y')\\
  &{}\cdot p'_{\alpha'+3\beta'}(-x'^3y')
  p'_{2\alpha'+3\beta'}(-x'^3y'^2),\\
 p'_{\alpha'+\beta'}(y')p'_{\beta'}(x')
  ={}&p'_{\beta'}(x')p'_{\alpha'+\beta'}(y')
  p'_{\alpha'+2\beta'}(x'y'),\\
 p'_{\alpha'+3\beta'}(y')p'_{\alpha'}(x')
  ={}&p'_{\alpha'}(x')p'_{\alpha'+3\beta'}(y')
  p'_{2\alpha'+3\beta'}(x'y').
\end{aligned}
\]
Define \(\phi:U\to U'\) by
\[
\begin{aligned}
 &\phi\!\left(
 p_\alpha(x)p_\beta(y)p_{\alpha+\beta}(t)p_{2\alpha+\beta}(u)
 p_{3\alpha+\beta}(v)p_{3\alpha+2\beta}(w)\right)\\
 &\quad=
 p'_{\alpha'}(x^{3q})p'_{\beta'}(y^q)
 p'_{\alpha'+3\beta'}(t^{3q})p'_{2\alpha'+3\beta'}(u^{3q})
 p'_{\alpha'+\beta'}(v^q)p'_{\alpha'+2\beta'}(w^q).
\end{aligned}
\]
Substituting \(x^{3q},y^q\) in the three displayed relations for
\(U'\) yields exactly the images of the three relations for \(U\);
thus \(\phi\) is a homomorphism.

\SGARefTarget{sga3:XXIII:4.1.9}
\SGAParagraph{4.1.9. Semisimple rank greater than \(2\)}
For \(\alpha\in\Delta\), put \(\alpha'=d(\alpha)\). For every
\((\alpha,\beta)\), consider the pinned rank-at-most-two groups
\[
  Z_{\alpha\beta},\qquad Z'_{\alpha'\beta'}.
\]
The homomorphism \(M'\to M\) underlying \(h\) induces a
\(p\)-morphism
\[
  h_{\alpha\beta}:R(Z'_{\alpha'\beta'})
    \longrightarrow R(Z_{\alpha\beta}).
\]
By the preceding cases, there is a morphism of pinned groups
\[
  f_{\alpha\beta}:Z_{\alpha\beta}
    \longrightarrow Z'_{\alpha'\beta'}
\]
with \(R(f_{\alpha\beta})=h_{\alpha\beta}\). The restrictions
\[
  f_{\alpha\beta}|_{Z_\alpha},
  \qquad f_{\alpha\alpha}:Z_\alpha\to Z'_{\alpha'}
\]
correspond to the same morphism of pinned root data, and therefore
coincide by uniqueness. \SGARefLink{sga3:XXIII:2.4}{Corollary 2.4} now glues the
\(f_{\alpha\beta}\) to a morphism \(f:G\to G'\). It is a morphism of
pinned groups and satisfies \(R(f)=h\), completing the proof.

% END inputs/XXIII/tex/components/09_expose_XXIII_fundamental_theorem.tex

% BEGIN inputs/XXIII/tex/components/10_expose_XXIII_corollaries.tex
\SGANextSectionTarget{sga3:xxiii:section:5}\SGASection{5}{Corollaries of the Fundamental Theorem}

The most important is the following.

\SGARefTarget{sga3:XXIII:5.1}
\SGAStatement{Corollary 5.1}{Let \(S\) be a nonempty scheme, let
\(G,G'\) be pinned \(S\)-groups, and let
\[
  h:R(G')\xrightarrow{\sim}R(G)
\]
be an isomorphism of pinned root data. There is a unique isomorphism
of pinned \(S\)-groups \(f:G\xrightarrow{\sim}G'\) such that
\(R(f)=h\).}

This also follows from \SGARefLink{sga3:XXIII:3.5.1}{Theorem 3.5.1}, because its relations can be
written using only \(R(G)\). It also follows from the elementary part
of the proof of \SGARefLink{sga3:xxiii:theorem:4-1}{Theorem 4.1}, without the exceptional isogenies of
\SGARefLink{sga3:XXIII:4.1.7}{4.1.7} and \SGARefLink{sga3:XXIII:4.1.8}{4.1.8}.

\SGARefTarget{sga3:XXIII:5.2}
\SGAStatement{Corollary 5.2 (Uniqueness theorem)}{Let \(S\) be a
scheme and let \(G,G'\) be two splittable \(S\)-groups. If they have
the same type, then they are isomorphic.}

\SGARefTarget{sga3:XXIII:5.3}
\SGAStatement{Corollary 5.3}{Let \(S\) be a scheme and \(G,G'\) two
splittable \(S\)-groups. The following are equivalent:
\begin{enumerate}
  \item[\SGAInlineRefTarget{sga3:XXIII:5.3:item:i}(i)] \(G\) and \(G'\) are isomorphic;
  \item[\SGAInlineRefTarget{sga3:XXIII:5.3:item:ii}(ii)] they are locally isomorphic for the fpqc topology;
  \item[\SGAInlineRefTarget{sga3:XXIII:5.3:item:iii}(iii)] there is \(s\in S\) such that \(G_s\) and \(G'_s\)
    have the same type.
\end{enumerate}}

Clearly \SGARefLink{sga3:XXIII:5.3:item:i}{(i)} implies \SGARefLink{sga3:XXIII:5.3:item:ii}{(ii)}, which implies \SGARefLink{sga3:XXIII:5.3:item:iii}{(iii)}. Conversely, \SGARefLink{sga3:XXIII:5.3:item:iii}{(iii)}
gives
\[
  R(G)=R(G_s)=R(G'_s)=R(G'),
\]
so \SGARefLink{sga3:XXIII:5.2}{Corollary 5.2} applies.

\SGARefTarget{sga3:XXIII:5.4}
\SGAStatement{Corollary 5.4 (Uniqueness of Chevalley schemes)}{Let
\(G,G'\) be reductive groups over \(\mathbf Z\), each possessing a
split maximal torus. The following are equivalent:
\begin{enumerate}
  \item[(i)] \(G\) and \(G'\) are isomorphic;
  \item[(ii)] there is \(s\in\operatorname{Spec}\mathbf Z\) such
    that \(G_s\) and \(G'_s\) have the same type;
  \item[(iii)] \(G_{\mathbf C}\) and \(G'_{\mathbf C}\) have the same
    type.
\end{enumerate}}

Indeed, \(G\) and \(G'\) are splittable by \SGARefLink{sga3:xxii:prop:2.2}{Exposé XXII, 2.2}.
\NDE{In fact every \(\mathbf Z\)-torus is split: every
\(\mathbf Z\)-torus is isotrivial, and every étale cover of
\(\operatorname{Spec}\mathbf Z\) is trivial.}

\SGARefTarget{sga3:XXIII:5.5}
\SGAStatement{Corollary 5.5 (Existence of outer automorphisms)}{Let
\(S\) be a scheme, \(G\) a pinned \(S\)-group, and \(h\) an
automorphism of the pinned root datum \(R(G)\). There is a unique
automorphism \(u\) of \(G\) preserving the pinning and satisfying
\(R(u)=h\).}

\SGARefTarget{sga3:XXIII:5.5bis}
\SGAStatement{Corollary 5.5 bis}{Let \(S\) be a scheme, \(G\) a
split reductive \(S\)-group, and \(R^+\) a system of positive roots.
For every simple root \(\alpha\), choose an isomorphism of vector
groups
\[
  p_\alpha:\mathbf G_{a,S}\xrightarrow{\sim}U_\alpha.
\]
Let \(h\) be an automorphism of \(M\) permuting the positive roots
and the corresponding coroots: if \(\alpha\in R^+\), then
\[
  h(\alpha)\in R^+,\qquad h^\vee(\alpha^*)=h(\alpha)^*.
\]
There is a unique automorphism \(u\) of \(G\) inducing \(D_S(h)\) on
\(T\) and satisfying \(u\circ p_\alpha=p_{h(\alpha)}\) for every
simple root \(\alpha\).}

\SGARefTarget{sga3:XXIII:5.6}
\SGAStatement{Corollary 5.6}{Let \(G,G'\) be reductive
\(S\)-groups. The following are equivalent:
\begin{enumerate}
  \item[\SGAInlineRefTarget{sga3:XXIII:5.6:item:i}(i)] \(G,G'\) are locally isomorphic for the fpqc topology;
  \item[(i bis)] they are locally isomorphic for the étale topology;
  \item[(ii)] \(G_s,G'_s\) are isomorphic for every \(s\in S\);
  \item[\SGAInlineRefTarget{sga3:XXIII:5.6:item:iii}(iii)] the functions \(s\mapsto\operatorname{type}(G_s)\)
    and \(s\mapsto\operatorname{type}(G'_s)\) are equal.
\end{enumerate}}

The implication (i bis)\(\Rightarrow\)\SGARefLink{sga3:XXIII:5.6:item:i}{(i)} is immediate;
(i)\(\Rightarrow\)(ii) follows from the principle of finite
extension (EGA IV\(_3\), 9.1.4); and
(ii)\(\Rightarrow\)(iii) is immediate. For
\SGARefLink{sga3:XXIII:5.6:item:iii}{(iii)}\(\Rightarrow\)(i bis), one may suppose \(G,G'\) splittable
(\SGARefLink{sga3:xxii:cor:2.3}{Exposé XXII, 2.3}), and then apply \SGARefLink{sga3:XXIII:5.3}{Corollary 5.3}.

\SGARefTarget{sga3:XXIII:5.7}
\SGAStatement{Corollary 5.7}{Let \(S\) be a scheme, \(G\) a
reductive \(S\)-group, and \(G'\) an affine smooth \(S\)-group with
connected fibers. If \(s\in S\) and \(G_s\simeq G'_s\), then there
is an étale morphism \(S'\to S\), covering \(s\), such that
\(G_{S'}\simeq G'_{S'}\).}

By Exposés XIX, \SGARefLink{sga3:XXIII:2.5}{2.5} and XXII, \SGARefLink{sga3:xxii:prop:2.1}{2.1}, one may suppose \(G,G'\)
reductive and splittable, and \SGARefLink{sga3:XXIII:5.3}{Corollary 5.3} applies.

\SGARefTarget{sga3:XXIII:5.8}
\SGAStatement{Corollary 5.8}{Let \(k\) be a field and \(G,G'\)
reductive \(k\)-groups. The following are equivalent:
\begin{enumerate}
  \item[(i)] \(G,G'\) have the same type;
  \item[(ii)] \(G\otimes_k\bar k\) and
    \(G'\otimes_k\bar k\) are isomorphic;
  \item[(iii)] for some finite separable extension \(K/k\),
    \(G\otimes_kK\) and \(G'\otimes_kK\) are isomorphic.
\end{enumerate}}

\SGARefTarget{sga3:XXIII:5.9}
\SGAStatement{Corollary 5.9}{Let \(S\) be a nonempty scheme and
\(R\) a root datum. The following are equivalent:
\begin{enumerate}
  \item[\SGAInlineRefTarget{sga3:XXIII:5.9:item:i}(i)] there is a pinned \(S\)-group of type \(R\);
  \item[(ii)] there is an \(S\)-group of type \(R\);
  \item[\SGAInlineRefTarget{sga3:XXIII:5.9:item:iii}(iii)] fpqc-locally, there is a reductive \(S\)-group of
    type \(R\).
\end{enumerate}}

Only \SGARefLink{sga3:XXIII:5.9:item:iii}{(iii)}\(\Rightarrow\)\SGARefLink{sga3:XXIII:5.9:item:i}{(i)} needs proof. Suppose, for simplicity,
that there are a faithfully flat quasi-compact morphism \(S'\to S\)
and a reductive \(S'\)-group \(G'\) of type \(R\). We may suppose
\(G'\) splittable and choose a pinning \(E'\), writing
\(R=R(G',E')\). The two pullbacks of \((G',E')\) to
\[
  S''=S'\times_SS'
\]
are pinned groups \((G''_1,E''_1)\), \((G''_2,E''_2)\), with
canonical isomorphisms
\[
  p_i:R(G''_i,E''_i)\xrightarrow{\sim}R.
\]
Thus
\[
  p=p_2^{-1}\circ p_1:
  R(G''_1,E''_1)\xrightarrow{\sim}R(G''_2,E''_2).
\]
The uniqueness theorem gives a unique isomorphism
\[
  f:(G''_1,E''_1)\xrightarrow{\sim}(G''_2,E''_2)
\]
with \(R(f)=p\). This is a descent datum on \(G'\). Its cocycle
condition may be checked after applying \(R\), since \(R\) is fully
faithful, and then follows from the construction of \(p\). Because
\(G'\) is affine, the datum is effective. The pinning is stable
under descent, so it descends to a pinned \(S\)-group \((G,E)\) of
type \(R\).

\medskip
\noindent\textit{N.B.}\quad
In the language of fibered categories, the preceding proof becomes
both shorter and more transparent.

\SGARefTarget{sga3:XXIII:5.10}
\SGAStatement{Corollary 5.10}{Let \(S\) be a nonempty scheme and
\(R\) a pinned root datum for which there exists a reductive
\(S\)-group of type \(R\). Then there is a pinned \(S\)-group of
type \(R\), unique up to unique isomorphism.}

\SGARefTarget{sga3:XXIII:5.11}
\SGAStatement{Definition 5.11}{Under these hypotheses, denote the
unique pinned \(S\)-group of type \(R\) by
\(\operatorname{\acute Ep}_S(R)\), its canonical maximal torus by
\(T_S(R)\), its canonical Borel subgroup by \(B_S(R)\), and so on.}

For every morphism \(S'\to S\), with \(S'\ne\varnothing\), one may
identify
\[
  \operatorname{\acute Ep}_{S'}(R)
    =\operatorname{\acute Ep}_S(R)\times_SS'.
\]
If \(\operatorname{\acute Ep}_{\operatorname{Spec}\mathbf Z}(R)\)
exists, write it \(\operatorname{\acute Ep}(R)\); then
\[
  \operatorname{\acute Ep}_S(R)
    =\operatorname{\acute Ep}(R)
      \times_{\operatorname{Spec}\mathbf Z}S.
\]
The group \(\operatorname{\acute Ep}(R)\) is called the
\emph{Chevalley group scheme of type \(R\)}.

\SGARefTarget{sga3:XXIII:5.12}
\SGAParagraph{5.12}
Thus an \(S\)-sheaf in groups \(G\) is a reductive \(S\)-group of
type \(R\) if and only if it is locally isomorphic, for the étale or
fpqc topology, to \(\operatorname{\acute Ep}_S(R)\). Similarly, by
the conjugacy theorems, giving a reductive \(S\)-group of type \(R\)
with a maximal torus is locally equivalent to
\[
  (\operatorname{\acute Ep}_S(R),T_S(R)).
\]
The analogous statements hold for Borel subgroups and Killing pairs.

% END inputs/XXIII/tex/components/10_expose_XXIII_corollaries.tex

% BEGIN inputs/XXIII/tex/components/11_expose_XXIII_chevalley_systems_and_bibliography.tex
\SGANextSectionTarget{sga3:xxiii:section:6}\SGASection{6}{Chevalley Systems}

The explicit calculations of \SGARefLink{sga3:xxiii:section:3}{Section 3} have important numerical
consequences.

\SGARefTarget{sga3:XXIII:6.1}
\SGAStatement{Definition 6.1}{Let \(S\) be a scheme and
\((G,T,M,R)\) a split \(S\)-group. A \emph{Chevalley system} of
\(G\) is a family
\[
  (X_\alpha)_{\alpha\in R},
  \qquad
  X_\alpha\in\Gamma(S,\mathfrak g^\alpha)^\times,
\]
satisfying:
\begin{description}
  \item[(SC)] for every pair \(\alpha,\beta\in R\),
  \[
    \operatorname{Ad}(w_\alpha(X_\alpha))X_\beta
      =\pm X_{s_\alpha(\beta)}.
  \]
\end{description}}

Recall that
\[
  w_\alpha(X_\alpha)
    =\exp(X_\alpha)\exp(-X_\alpha^{-1})\exp(X_\alpha).
\]
Condition (SC) implies, in particular,
\[
  X_{-\alpha}=\pm X_\alpha^{-1}
\]
by the identity
\[
  \operatorname{Ad}(w_\alpha(X_\alpha))X_\alpha=-X_\alpha^{-1}.
\]

\SGARefTarget{sga3:XXIII:6.2}
\SGAStatement{Proposition 6.2}{Every split group admits a Chevalley
system. More precisely, if
\[
  \bigl(\Delta,(X'_\alpha)_{\alpha\in\Delta}\bigr)
\]
is a pinning of \((G,T,M,R)\), there is a Chevalley system
\((X_\alpha)_{\alpha\in R}\) with \(X_\alpha=X'_\alpha\) for
\(\alpha\in\Delta\).}

First, it suffices to verify (SC) for \(\alpha\in\Delta\). Indeed,
for every \(\alpha\in R\), choose a sequence
\(\alpha_i\in\Delta\) such that
\[
  \alpha=s_{\alpha_1}\cdots s_{\alpha_n}(\alpha_{n+1}).
\]
Repeated use of (SC) gives
\[
 X_\alpha
  =\pm\operatorname{Ad}\!\left(
    w_{\alpha_1}(X_{\alpha_1})\cdots
    w_{\alpha_n}(X_{\alpha_n})\right)X_{\alpha_{n+1}}.
\]
By \SGARefLink{sga3:xxiii:lemma:3-1-1:item:iii}{Lemma 3.1.1(iii)},
\[
\begin{aligned}
 &w_{\alpha_1}(X_{\alpha_1})\cdots w_{\alpha_n}(X_{\alpha_n})
  w_{\alpha_{n+1}}(X_{\alpha_{n+1}})
  w_{\alpha_n}(X_{\alpha_n})^{-1}\cdots
  w_{\alpha_1}(X_{\alpha_1})^{-1}\\
 &\hspace{42mm}=\alpha^*(\pm1)w_\alpha(X_\alpha).
\end{aligned}
\]
Now
\[
 w_{\alpha_i}(X_{\alpha_i})^{-1}
  =\alpha_i^*(-1)w_{\alpha_i}(X_{\alpha_i}),
\]
and
\[
 \beta(\gamma^*(-1))=(-1)^{(\gamma^*,\beta)}=\pm1.
\]
It follows that (SC) for every pair \((\alpha_i,\gamma)\) implies
(SC) for every \((\alpha,\beta)\).

For each \(\alpha\in R\), choose such a sequence and define
\[
 X_\alpha=
 \operatorname{Ad}\!\left(
  w_{\alpha_1}(X_{\alpha_1})\cdots
  w_{\alpha_n}(X_{\alpha_n})\right)X'_{\alpha_{n+1}}.
\]
It remains to prove the following lemma.

\SGARefTarget{sga3:XXIII:6.3}
\SGAStatement{Lemma 6.3}{Let
\[
 (G,T,M,R,\Delta,(X_\alpha)_{\alpha\in\Delta})
\]
be a pinned \(S\)-group. If \(\alpha_i\), \(0\le i\le n+1\), are
simple roots such that
\[
  s_{\alpha_1}\cdots s_{\alpha_n}(\alpha_{n+1})=\alpha_0,
\]
then
\[
  \operatorname{Ad}(w_{\alpha_1}\cdots w_{\alpha_n})
    X_{\alpha_{n+1}}=\pm X_{\alpha_0}.
\]}

Arguing as in \SGARefLink{sga3:XXIII:2.3.4}{2.3.4}, using Tits's lemma, reduces the proof to:
\begin{enumerate}
  \item[\SGAInlineRefTarget{sga3:XXIII:6.3:item:a}(a)] \(G\) has semisimple rank at most \(2\); or
  \item[\SGAInlineRefTarget{sga3:XXIII:6.3:item:b}(b)] \(w_{\alpha_1}\cdots w_{\alpha_n}\) is a section of
    \(T\).
\end{enumerate}
Case \SGARefLink{sga3:XXIII:6.3:item:a}{(a)} is exactly the calculation in \SGARefLink{sga3:XXIII:3.1.2:item:i}{3.1.2(i)}, \SGARefLink{sga3:xxiii:proposition:3-2-1:item:i}{3.2.1(i)},
\SGARefLink{sga3:xxiii:proposition:3-3-1:item:i}{3.3.1(i)}, or \SGARefLink{sga3:xxiii:proposition:3-4-1:item:i}{3.4.1(i)}.

In case \SGARefLink{sga3:XXIII:6.3:item:b}{(b)}, it suffices to show that if
\[
  s_{\alpha_1}\cdots s_{\alpha_n}=1,
\]
then \(t=w_{\alpha_1}\cdots w_{\alpha_n}\) satisfies
\(\alpha(t)=\pm1\) for every root \(\alpha\). By the structure of the
Weyl group, the word in the free group generated by the
\(s_\alpha\) lies in the normal subgroup generated by
\[
  (s_\alpha s_\beta)^{n_{\alpha\beta}}.
\]
Thus one reduces to a word
\[
 s_{\alpha_1}\cdots s_{\alpha_i}
 (s_{\alpha_{i+1}}s_{\alpha_{i+2}})
 ^{n_{\alpha_{i+1}\alpha_{i+2}}}
 s_{\alpha_i}\cdots s_{\alpha_1},
\]
for which
\[
  t=s_{\alpha_1}\cdots s_{\alpha_i}
      (t_{\alpha_{i+1}\alpha_{i+2}}).
\]
The assertion is now immediate from the values of
\(t_{\alpha_1\alpha_2}\) computed in \SGARefLink{sga3:XXIII:3.1.2:item:i}{3.1.2(i)}, \SGARefLink{sga3:xxiii:proposition:3-2-1:item:i}{3.2.1(i)},
\SGARefLink{sga3:xxiii:proposition:3-3-1:item:i}{3.3.1(i)}, and \SGARefLink{sga3:xxiii:proposition:3-4-1:item:i}{3.4.1(i)}.

\SGARefTarget{sga3:XXIII:6.4}
\SGAStatement{Proposition 6.4}{Let \((G,T,M,R)\) be a split
\(S\)-group with Chevalley system \((X_\alpha)_{\alpha\in R}\).
Let \(\alpha,\beta\) be nonproportional roots, with
\[
  \operatorname{length}(\alpha)\le\operatorname{length}(\beta),
  \qquad
  |(\beta^*,\alpha)|\le|(\alpha^*,\beta)|.
\]
Let \(q\) and \(p-1\) be the nonnegative integers for which the roots
of the form \(\beta+k\alpha\) are
\[
  \{\beta-(p-1)\alpha,\ldots,\beta,\ldots,\beta+q\alpha\}.
\]
Thus \(-(\alpha^*,\beta)=q-p+1\). Writing
\(p_\gamma(x)=\exp(xX_\gamma)\), the commutation relation between
\(U_\alpha\) and \(U_\beta\) is given by:
\[
\begin{array}{c|l}
(p,q)&p_\beta(y)p_\alpha(x)p_\beta(-y)p_\alpha(-x)\\ \hline
(-,0)&e\\
(1,1)&p_{\alpha+\beta}(\pm xy)\\
(1,2)&p_{\alpha+\beta}(\pm xy)p_{2\alpha+\beta}(\pm x^2y)\\
(1,3)&p_{\alpha+\beta}(\pm xy)p_{2\alpha+\beta}(\pm x^2y)
       p_{3\alpha+\beta}(\pm x^3y)p_{3\alpha+2\beta}(\pm x^3y^2)\\
(2,1)&p_{\alpha+\beta}(\pm2xy)\\
(2,2)&p_{\alpha+\beta}(\pm2xy)p_{2\alpha+\beta}(\pm3x^2y)
       p_{\alpha+2\beta}(\pm3xy^2)\\
(3,1)&p_{\alpha+\beta}(\pm3xy)
\end{array}
\]
These are all possibilities, since the root string has length at
most \(3\).}

\SGAProof
By \SGARefLink{sga3:XXI:3.5.4}{Exposé XXI, 3.5.4}, there is a system of simple roots \(\Delta\)
such that \(\alpha\in\Delta\), and there are
\(\alpha'\in\Delta\), \(a,b\in\mathbf Q_+\), with
\[
  \beta=a\alpha+b\alpha'.
\]
The desired relation lies in \(U_{\alpha\alpha'}\), so it follows
from the rank-two calculations of Section 3 and condition (SC).

\SGARefTarget{sga3:XXIII:6.5}
\SGAStatement{Corollary 6.5 (Chevalley's rule)}{Let
\((X_\alpha)_{\alpha\in R}\) be a Chevalley system of the split
\(S\)-group \(G\). If \(\alpha,\beta,\alpha+\beta\in R\), then
\[
  [X_\alpha,X_\beta]=\pm pX_{\alpha+\beta},
\]
where \(p\) is the least positive integer for which
\(\beta-p\alpha\) is not a root.}

The assertion is symmetric in \(\alpha,\beta\); after arranging
\(\operatorname{length}(\alpha)\le\operatorname{length}(\beta)\),
it follows from \SGARefLink{sga3:XXIII:6.4}{Proposition 6.4}.

\SGARefTarget{sga3:XXIII:6.6}
\SGAStatement{Corollary 6.6}{Let \(S\) be a scheme such that
\(6\cdot1_S\ne0\), and let \((G,T,M,R)\) be a split \(S\)-group.
If \(R'\subset R\) is such that
\begin{equation}
  \mathfrak g_{R'}
   =\mathfrak t\oplus\bigoplus_{\alpha\in R'}\mathfrak g^\alpha
  \tag{*}\label{sga3:XXIII:6.6:star}
\end{equation}
is a Lie subalgebra of \(\mathfrak g\), then \(R'\) is a closed
subset of \(R\).}

Choose a Chevalley system and let \(\alpha,\beta\in R'\) with
\(\alpha+\beta\in R\). By \SGARefLink{sga3:XXIII:6.5}{Corollary 6.5} and \SGARefLink{sga3:XXIII:6.4}{Proposition 6.4},
\[
  [X_\alpha,X_\beta]=\pm pX_{\alpha+\beta},
  \qquad p\in\{1,2,3\}.
\]
Since neither \(2\) nor \(3\) is zero on \(S\),
\eqref{sga3:XXIII:6.6:star} implies
\(\mathfrak g^{\alpha+\beta}\subset\mathfrak g_{R'}\), and hence
\(\alpha+\beta\in R'\).
\NDE{It suffices that \(2\) and \(3\) are nonzero on \(S\); for
example the result applies to \(S=\operatorname{Spec}(\mathbf Z/6\mathbf Z)\).
In the excluded characteristics, short-root subsystems furnish
standard counterexamples.}

\SGARefTarget{sga3:XXIII:6.7}
\SGAParagraph{6.7}
The exact signs in this section can be determined by studying
\(\operatorname{Norm}_G(T)\), and more precisely the extended Weyl
group
\[
  W^*=N(\mathbf Z),\qquad
  N=\operatorname{Norm}_{\operatorname{\acute{E}p}_{\mathbf Z}(R)}
         (T_{\mathbf Z}(R)).
\]
It is an extension of \(W(R)\) by an abelian group of type
\((2,2,\ldots,2)\), which accounts for the signs.

\SGARefTarget{sga3:XXIII:6.7.1}
\SGAStatement{Remark 6.7.1}{By 3.1.2(i) and 3.\(n\).1,
\(n=2,3,4\), the elements \(w_\alpha,w_\beta\in N(\mathbf Z)\)
satisfy the braid relations
\[
  w_\alpha w_\beta\cdots=w_\beta w_\alpha\cdots
\]
with \(n_{\alpha\beta}\) factors on each side.}

See also Tits \cite{Ti66}, Bourbaki \cite{BLie}, IX.4, Ex.~12, and
Springer \cite[9.3.2]{Sp98}.

\begin{thebibliography}{99}
\bibitem[BLie]{BLie}
N.~Bourbaki,
\emph{Groupes et algèbres de Lie}, Chap.~IX,
Masson, 1982.

\bibitem[Ja87]{Ja87}
J.~C. Jantzen,
\emph{Representations of Algebraic Groups},
Academic Press, 1987; 2nd ed., American Mathematical Society, 2003.

\bibitem[Sp98]{Sp98}
T.~A. Springer,
\emph{Linear Algebraic Groups}, 2nd ed.,
Birkhäuser, 1998.

\bibitem[St62]{St62}
R.~Steinberg,
\emph{Générateurs, relations et revêtements de groupes algébriques},
Colloque sur la théorie des groupes algébriques (Bruxelles, 1962),
Université de Louvain and Gauthier--Villars, 1962;
pp.~133--147 in \emph{R. Steinberg Collected Papers},
American Mathematical Society, 1997.

\bibitem[St67]{St67}
R.~Steinberg,
\emph{Lectures on Chevalley Groups},
Yale University, 1967.

\bibitem[Ta83]{Ta83}
M.~Takeuchi,
\emph{A hyperalgebraic proof of the isomorphism and isogeny
theorems for reductive groups},
J. Algebra \textbf{85} (1983), 179--196.

\bibitem[Ti66]{Ti66}
J.~Tits,
\emph{Normalisateurs de tores I. Groupes de Coxeter étendus},
J. Algebra \textbf{4} (1966), 96--116.
\end{thebibliography}

% END inputs/XXIII/tex/components/11_expose_XXIII_chevalley_systems_and_bibliography.tex


